%% Advances in Agriculture (Wiley) submission -- Manuscript 1
%% Built on the Wiley LaTeX authoring template (USG.cls, Wiley PDF design).
%% Compile: pdflatex -> bibtex -> pdflatex x2   (Times fallback per template guidelines)
\documentclass[ASNA,twocolumn]{USG}

\usepackage[T1]{fontenc} % single-glyph accented letters; correct extraction
\renewcommand{\rmdefault}{ptm}% Times text (guideline fallback for pdflatex); math stays in the class setup
\graphicspath{{./}{./images/}}

% Referenceable anchors for the unnumbered appendix headings produced by \bmsection.
\makeatletter
\newcommand{\applabel}[2]{\phantomsection\edef\@currentlabel{#1}\label{#2}}
\makeatother

\articletype{Research Article}%
% History dates are publisher-assigned. The Wiley author guidelines instruct authors to
% enter "00" while the manuscript is being prepared; empty braces print bare labels.
\received{00}
\revised{00}
\accepted{00}
\journal{Advances in Agriculture}
% \copyyear is deliberately unset: it drives a footer asserting a publication year and
% page range, both of which the publisher assigns.
\startpage{1}
\articledoi{}

\begin{document}

\title{Global Geography of AI-Based Crop-Yield Research: Agricultural Scale, Research Need, and Scientific Capacity}

\author[1]{Prasamsha Poudel}[https://orcid.org/0009-0003-6441-3023]
\author[2]{Saugat Khanal}[https://orcid.org/0000-0002-9146-7428]

\authormark{POUDEL and KHANAL}
\titlemark{Global Geography of AI-Based Crop-Yield Research}

\address[1]{\orgname{Agriculture and Forestry University, }%
\orgaddress{\city{Puranchaur, Pokhara, }\state{Kaski 33713, }\country{Nepal}}}

\address[2]{\orgdiv{Department of Agricultural Economics, }\orgname{Oklahoma State University, }%
\orgaddress{\city{Stillwater, }\state{Oklahoma 74078, }\country{USA}}}

\corres{Prasamsha Poudel (\email{poudelprasamsha64@gmail.com}) ~|~ Saugat Khanal (\email{saugat.khanal@okstate.edu})}

\keywords{crop-yield prediction | artificial intelligence | machine learning | research allocation | scientific capacity | agricultural R\&D | spatial inequality}

\abstract[ABSTRACT]{AI-based crop-yield prediction has expanded quickly, but the evidence it produces need
not come from the places where crops occupy the most land or where additional evidence would be
most useful. This study maps that geography against agricultural scale, research need, and
scientific capacity. From a frozen
corpus of 7,045 eligible studies published between 2000 and 2026, each study's empirical location was read from its reported study area and data rather than from
author affiliations, then allocated fractionally over 2,616 country-crop systems covering 199
countries and 25 crops. Evidence is concentrated: 122 countries carry at least one eligible study, 485 systems carry
research, and the Gini coefficient across the 199 countries is 0.848. Harvested area has the largest association with
whether a system is studied: an interquartile-range increase raises the predicted probability by
35 percentage points, about five times the change from any other covariate, and the association
holds under crop fixed effects and within each of the five most studied crops.
Research and development expenditure is positively associated with both participation and
intensity; the remaining capacity indicators do not point one way. The association with measured research need is not distinguishable from zero for participation.
A positive intensity association appears only when research is constrained to scale
proportionally with harvested area; estimating that scaling freely eliminates it. Research standing relative to need is spatially clustered (Moran's $I = 0.363$), with
High--High clusters in western and northern Europe and Low--Low clusters in central and eastern
Africa.}
% Abbreviations are defined at first use in the body. \abbr is a class no-op and
% would silently discard anything passed to it, so no list is supplied here.

\contributed{Prasamsha Poudel and Saugat Khanal contributed equally to this study.}

% Copyright/licence line is applied by the publisher at production; none is printed on the submission.

\maketitle

\section{Introduction}\label{sec:intro}


Reliable crop-yield information supports decisions made before and during the growing season,
including input allocation, production planning, and responses to weather and market
uncertainty. At broader scales, yield estimates also inform food-supply assessments,
agricultural policy, and resource management \cite{vanKlompenburg2020}. Improving the accuracy
and practical value of yield prediction remains an important empirical problem in
agricultural research.

Recent reviews document rapid growth in AI- and machine-learning-based crop-yield prediction,
with models combining weather, soil, vegetation, and remote-sensing data across a wide range of
statistical and machine-learning approaches \cite{Shawon2025,vanKlompenburg2020}. The expansion
of AI-based yield prediction does not establish that more complex methods uniformly outperform
simpler alternatives. Reported accuracy depends on the available data, crop and production setting,
and evaluation design \cite{Jabed2024}.

Geographic transferability is therefore central to the interpretation of crop-yield models.
High within-sample accuracy does not guarantee similar performance in new fields, regions, or
production systems. Studies using spatially independent validation generally report weaker
out-of-domain performance than random validation, indicating that spatial structure can
inflate apparent predictive accuracy \cite{Habibi2024,Okupska2026,Stiller2024}. The geographic scope of the evidence base affects how broadly reported model
performance can be interpreted.

The geographic distribution of crop-yield research also reflects a broader pattern in
scientific production. Cross-country studies show that research output remains concentrated
in countries with stronger scientific institutions and research resources, although
participation has broadened over time \cite{Bueno2025,CashGibson2018}. Persistent
concentration can shape both the problems that receive research attention and the locations
where empirical evidence is generated.

The agricultural R\&D literature introduces a related distinction among agricultural importance,
research need, and scientific capacity. Agricultural importance reflects the scale of crop
production, measured here by agricultural scale, primarily harvested area; research need reflects problems such as food insecurity, low productivity,
vulnerability, or yield gaps; and scientific capacity reflects the resources and institutions
available to conduct research. Those dimensions often diverge geographically
\cite{NinPratt2021,NinPratt2023,Rosegrant2023,vanDijk2025}.

Publication counts alone provide an incomplete measure of research allocation. A
country may account for many crop-yield studies because it has a large scientific base,
because it is a major producer, or both. Publication totals cannot distinguish research
concentration that is proportional to agricultural scale from concentration that
reflects unequal scientific capacity. The same problem arises when research attention is
compared with research need.

This study addresses that gap by mapping AI- and machine-learning-based crop-yield research
using empirical study locations rather than author affiliations. Research attention is
evaluated at the country-crop level against production and harvested area, broader indicators
of research need, and national scientific capacity. The country-crop design allows research
allocation to be compared separately with agricultural scale, research need, and
scientific capacity, while also permitting direct assessment of the spatial pattern of
resulting mismatches.

\section{Literature Review}

\subsection{AI- and Machine-Learning-Based Crop-Yield Prediction}

Reviews of AI-based yield prediction converge on the predictors the field relies on. Weather,
soil, and vegetation variables recur across studies, and remote sensing now supplies much of
the input, with satellite and UAV imagery feeding models that increasingly fuse several data
sources rather than depending on one \cite{vanKlompenburg2020,Joshi2023,Muruganantham2022,Shawon2025}.

No single model class dominates the recent literature. Earlier reviews found artificial neural
networks the most common approach \cite{vanKlompenburg2020}; later ones report
penalised regression, tree ensembles including gradient boosting, support vector machines,
and convolutional and recurrent architectures used side by side, together with hybrid and
stacked designs \cite{Jabed2024,Screpnik2025,Shawon2025,Sharma2023}.

That breadth has produced many reports of strong performance, but not a ranking in which
more complex models are reliably better. Reported accuracy varies with data quality, with
the crop and production setting, and with how evaluation is carried out
\cite{Screpnik2025,Shawon2025}. \citet{Jabed2024} treat yield estimation as a
problem shaped by geographical variation, crop diversity, and cultivation area, which makes
results from different studies difficult to compare.

A major limitation is therefore comparability across studies rather than the availability of
model classes. Reviews report
recurring problems with data quality, pre-processing, and model complexity
\cite{Jabed2024,Screpnik2025}, and although RMSE, $R^2$, and MAE are reported almost
everywhere, a shared metric is not a shared validation standard
\cite{Screpnik2025,Shawon2025}. Growth in this literature is evidence of methodological
expansion; it is not yet evidence that performance holds across crops and places.

\subsection{Spatial Transferability and Geographic Representation}

Strong performance within the training data does not guarantee reliable prediction
elsewhere. Extrapolation across fields and regions is difficult because yield relationships
are shaped by spatial autocorrelation, environmental heterogeneity, and processes that
operate differently at different scales \cite{Habibi2024,Stiller2024}.

Across studies that test this directly, randomly assigned validation folds give a more
optimistic view of transferability than geographically independent ones, in deep learning on
small fields \cite{Stiller2024}, in soybean prediction from UAV imagery
\cite{Habibi2024}, and in NDVI-based winter oilseed rape, where tree ensembles most
accurate under random sampling explained almost no variance under field-wise validation
\cite{Okupska2026}. These studies attribute the gap to models memorising spatial structure rather than learning
relationships that hold beyond the calibration domain.

Where these studies diverge is on which model classes retain predictive performance under
spatial extrapolation, and complexity does not help by default.
\citet{Habibi2024} found LASSO with recursive feature elimination transferring better than
richer models, while \citet{Okupska2026} found a deep network more stable than tree
ensembles under the same spatial split. Validation design determines what a reported accuracy
means; the ranking of model families does not survive a change of design. External validity is
therefore not a by-product of predictive accuracy, and evidence concentrated in a narrow set
of agroecological settings would overstate how broadly its conclusions apply.

\subsection{Scientific Capacity and Agricultural Research Allocation}

Research production is itself geographically uneven, and the bibliometric literature is
careful about what that unevenness measures. \citet{CashGibson2018} treat national output
as only a crude indication of research capacity, because it also reflects investment,
population, and institutional support, and they normalise productivity by population or GDP
per capita for that reason.

Concentration is also relational rather than merely distributive: \citet{Bueno2025} find
collaboration largely confined within a small core group rather than running between
core and periphery, so peripheral systems are not simply smaller but differently connected.
Neither study describes a fixed geography, because participation has broadened and capacity
outside the traditional core has grown \cite{CashGibson2018}.

In the agricultural R\&D literature, agricultural importance, research need, and scientific
capacity overlap imperfectly: a country may be agriculturally important but have a weak
research system, or a strong research system with little dependence on agriculture.

The policy weight attached to public agricultural R\&D is visible in the data infrastructure
built around it \cite{Fuglie2019}: \citet{vanDijk2025} assemble harmonized public
agricultural research spending and personnel for 190 countries from roughly 1960 to 2022,
because long-run series are needed to study productivity effects.
\citet{Rosegrant2023} estimate a central benefit-cost ratio of 33 for expanded
agricultural R\&D in the Global South, placing it among the highest-return development
investments, while noting that uncertainty in long-run estimation remains a legitimate
concern.

Measures of agricultural R\&D intensity complicate a simple interpretation of underinvestment.
\citet{NinPratt2021} argues that the conventional intensity ratio, R\&D spending as a
share of agricultural GDP, understates the effort developing countries make, yet finds that
low- and middle-income countries still face an investment gap close to 50\% of R\&D spending.
Low intensity ratios do not map mechanically onto low effort, but substantial constraints
remain. Those constraints track system size and institutional quality:
\citet{NinPratt2023} show that small agricultural research systems in low- and
middle-income countries tend to have weaker education and research systems, poorer innovation
environments, and poorer links to markets, and \citet{Stads2020} report Southeast Asian
research intensity falling from 0.50\% of agricultural GDP in 2000 to 0.33\% in 2017 as
spending stagnated against rising output.

Research need and agricultural investment are also imperfectly aligned. \citet{Zhao2023}
find food-insecure countries
receiving only 20\% of global agricultural investment, in fewer and smaller projects, and
\citet{Malec2024} show that the effects of agricultural innovation investment on food
security vary by subregion within sub-Saharan Africa, so need, impact, and capacity do not
align uniformly even within one world region.

\subsection{Existing Agricultural-AI Bibliometrics and Remaining Gap}

Bibliometric work on agricultural AI has mapped the field's growth and structure in detail.
For crop-yield prediction, these analyses report sustained publication growth and identify the
countries, institutions, and themes that dominate it
\cite{Lokeshwari2024,Momenpour2024}. \citet{Momenpour2024} additionally note that
food security in developing countries receives less attention than its importance would
suggest, which is a statement about coverage rather than volume.

The same concentration recurs across broader and more specific bibliometrics: collaboration
centred on a handful of countries in agricultural AI \cite{Bhagat2022,Zhang2021}, most
countries contributing fewer than 100 publications each in crop classification from remote
sensing \cite{SanchezChavez2026}, and a United States and China core in maize-yield
prediction with India and Brazil entering more recently \cite{Sakthivel2026}.
Satellite-based yield prediction shows the same geographic concentration alongside a shift
from regression toward learned models and sensor fusion \cite{Radocaj2026a}.

The principal geographic unit in these bibliometric studies is author affiliation,
institution, or collaboration network rather than the location of the agricultural system
examined. Even a crop-specific treatment such as the olive-yield literature is organised
around co-authorship structure and citation performance rather than around the production
systems the studies examine \cite{Toscano2026} (Table~\ref{tab:synthesis}).

\begin{table*}[!t]
\centering
\footnotesize
\setlength{\tabcolsep}{4pt}
\renewcommand{\arraystretch}{1.15}
\caption{How this study differs from related strands of literature}
\label{tab:synthesis}
\begin{tabular}{@{}>{\raggedright\arraybackslash}p{2.05cm}
                  >{\raggedright\arraybackslash}p{3.15cm}
                  >{\raggedright\arraybackslash}p{2.25cm}
                  >{\raggedright\arraybackslash}p{3.05cm}
                  >{\raggedright\arraybackslash}p{3.35cm}@{}}
\toprule
Literature strand & Primary research focus & Typical geographic unit & Typical outcome &
Main limitation relative to the present study \\
\midrule
AI/ML crop-yield prediction &
Predictor sets, model families, and reported accuracy &
Study site, field, or administrative region &
Predictive performance for a fitted crop and place &
Site selection is not the object of study, so coverage across crops and countries is
incidental \\
\addlinespace
Spatial transferability and external validation &
Whether fitted relationships hold outside the calibration domain &
Fields or spatial folds within one study area &
Performance loss under spatially aware validation &
Establishes that geography matters for validity without mapping where the evidence base
sits \\
\addlinespace
Scientific capacity and agricultural research allocation &
Research investment, capacity, and returns &
Country, region, or income group &
R\&D spending, intensity ratios, and productivity or welfare returns &
Measures money and capacity rather than published evidence, and rarely reaches the crop
level \\
\addlinespace
Agricultural-AI and crop-yield bibliometrics &
Publication growth, contributors, collaboration, and themes &
Author affiliation country or institution &
Output counts, co-authorship networks, and keyword structure &
Records where work is authored rather than where the agriculture studied is located, and
uses no agricultural denominator \\
\addlinespace
Present study &
Global distribution of AI/ML crop-yield evidence relative to agricultural scale,
research need, scientific capacity, and spatial concentration &
Country-crop study-location unit &
Research allocation and mismatch relative to production, harvested area, need, and
capacity &
--- \\
\bottomrule
\end{tabular}
\end{table*}

The literature identified in this review does not allocate empirical studies to the country-crop systems they examine while jointly
comparing research allocation with agricultural scale, research need, and scientific capacity. Doing so requires coding each
study's empirical location rather than its authors' affiliations, and evaluating the result
against agricultural and socioeconomic denominators at the same unit.


\section{Data}

The analysis draws on four bodies of data: a frozen set of eligible AI/ML crop-yield studies
coded to the country-crop systems each study examines, FAOSTAT production and harvested-area
series, country-level indicators of research need and scientific capacity, and a country
boundary layer. Table~\ref{tab:data} gives the source, unit, reference period, and analytical
role of each. Reproducibility detail supporting this section is in
Appendices~\ref{app:corpus}--\ref{app:variables}.

\subsection{AI/ML Crop-Yield Study Database}

The study database was built from OpenAlex. Retrieval returned 102,892 records, reduced to
43,543 unique works after duplicate resolution. All 43,543 were screened on title and
abstract against six eligibility criteria. An eligible study reports an empirical application in which the
yield of an agricultural crop is the modelled outcome, estimated with a data-driven
predictive method, in an eligible document type published from 2000 onward; process-based
crop simulation without a learned component does not qualify. Screening left a frozen corpus
of 7,045 eligible studies published between 2000 and 2026. \ref{app:corpus} reports
the retrieval design, the duplicate-resolution rules, the screening procedure, and the
measured screening error rates.

\subsection{Study-Location and Country-Crop Coding}

The unit of analysis is the country-crop system, so each study is assigned to the countries
and crops it empirically examines. Study location is coded from what a study reports about
the origin of its agricultural data, never from author affiliation. The two differ
systematically, and coding from affiliation would reproduce the geography of authorship
rather than measure the geography of evidence. Country names are accepted only where the
surrounding text asserts an empirical study area; names appearing in institutional,
bibliographic, or varietal contexts are rejected. Studies that are global or regional in scope, that name no study area, or that name only an
ambiguous place receive distinct scope labels. Where the coding accepted explicit
country evidence inside such a study, those countries enter the allocation; a study with no
accepted country evidence stays unallocated, so absence of a coded location remains separable
from absence of research.

Location and crop coding were carried out by an automated, model-assisted extractor rather than
by human reading, and every record is stored with its evidence text, a confidence level, and an
unverified status. Accuracy was estimated from a stratified probability sample of 210 records
adjudicated against the coded evidence. The countries that receive research weight are correct
for 90.4\% of adjudicated records (95\% CI 84.3 to 94.3), and 0.7\% are allocated wholly to the
wrong country. Recall is the weaker property: 62.0\% of sampled unresolved records (95\% CI 48.2
to 74.1) name a study area that the coder did not resolve. \ref{app:coding} reports the design,
the error taxonomy, and a check on whether the unresolved studies differ geographically.

Crops are standardised to a controlled vocabulary and matched to FAOSTAT items. Studies
covering several countries or crops are allocated fractionally, each carrying a total weight
of one divided equally across the combinations it examines, which preserves the study total
and prevents multi-country work from counting repeatedly. Full counts are retained for
comparison. Accepted country evidence and a FAOSTAT-matched crop coincide for 2,031 studies, which are
the studies that enter the allocation. \ref{app:coding} gives the coding rules, the cue
taxonomy, the resolution counts, and the full accounting from corpus to panel.

\subsection{Agricultural Scale}

Agricultural scale is the empirical proxy for agricultural importance, measured primarily by
harvested area and secondarily by production compared within a crop. Harvested area does not
measure economic value, food-security significance, or strategic priority, and no claim of that
kind is drawn from it. FAOSTAT production and harvested-area data provide the benchmarks.
Each country-crop cell takes the mean of the 2018--2022 reference period, a window balancing
currency against year-to-year variation in weather and reporting. Availability of production
and harvested-area values determines which cells can enter denominator-based comparisons.
\ref{app:variables} reports the extract, the identifier crosswalk, and the resulting cell
counts.

\subsection{Research Need and Scientific Capacity}

Need and capacity are measured separately so that the two can be tested against each other.
The research-need index combines nine country-level indicators spanning food security,
agricultural dependence, climate exposure, and production instability, each converted to a
percentile rank before aggregation. Two further FAOSTAT quantities, a country's share of
global harvested area and of global production, are recorded with the need data but excluded
from the index by design: they are the paper's agricultural-importance measures, and an index
containing them could not support a claim that importance and need are measured separately.
Rank aggregation is the primary index; equal weighting, entropy weighting, the first
principal component, and an eleven-component variant that adds the two scale shares back are
retained for sensitivity analysis.

Scientific capacity uses World Bank indicators of research and development effort, economic
resources, tertiary education, digital access, and population. Research and development
expenditure as a share of GDP is the primary capacity measure and the sparsest series, so its
coverage governs the estimation sample; GDP per capita, tertiary enrolment, internet use, and
population enter alongside it. Missing capacity values are not imputed.
\ref{app:variables} and Table~\ref{tab:vardetail} give the full indicator list, sources,
coverage, directions, transformations, and missingness.

\subsection{Spatial Data}

Country geometries come from the Natural Earth admin-0 boundaries, giving an analytical layer
of 195 countries. Four panel territories have no matching geometry and cannot enter the
spatial analysis; all four carry no eligible study, so their absence removes no research
observation.

\subsection{Analytical Sample}

The master panel is the set of country-crop cells for which FAOSTAT reports the crop in the
country at any year from 1990 to 2024, taken over the 199 countries and 25 crops carried through
the crosswalk, together with 30 cells for four territories that FAOSTAT reports inside another
country. That rule gives 2,616 cells, of which 485 carry at least one eligible study and 2,131
carry none. A further 109 cells carry no FAOSTAT area, production, or yield value for the reference period
and support no denominator-based comparison. None of the 109 carries research, and restricting
the panel to the 2,507 cells with a FAOSTAT observation raises the share of studied systems from
18.5\% to 19.3\%. The 2,131 empty cells record zero coded research attention: the
search covered every country, so a zero means no resolved eligible study was allocated to that
cell, not that no eligible study exists for it. Zeros stay distinct from cells where external
data are unavailable. Analytical
components use different subsets, and the differences are large enough that one sample size
would misrepresent them. Comparisons against agricultural scale use cells with a positive
production or harvested-area denominator; spatial analysis uses the 195 countries in the
boundary layer; the econometric models use the 1,799 cells across 142 countries for which
every covariate is observed, research and development expenditure being the covariate with the
most missing data. Per-covariate missingness is in Table~\ref{tab:vardetail}. No country is excluded on grounds of
research output. Table~\ref{tab:summstats} summarises the variables entering the models, and
Figure~\ref{fig:flow} traces the path from retrieved records to the estimation sample.

\begin{table*}[!t]
\centering
\footnotesize
\setlength{\tabcolsep}{4pt}
\renewcommand{\arraystretch}{1.15}
\caption{Data sources and measures used in the analysis}
\label{tab:data}
\begin{tabular}{@{}>{\raggedright\arraybackslash}p{2.15cm}
                  >{\raggedright\arraybackslash}p{3.05cm}
                  >{\raggedright\arraybackslash}p{2.45cm}
                  >{\raggedright\arraybackslash}p{1.95cm}
                  >{\raggedright\arraybackslash}p{1.75cm}
                  >{\raggedright\arraybackslash}p{2.75cm}@{}}
\toprule
Data component & Main measures & Source & Unit & Reference period & Analytical role \\
\midrule
AI/ML crop-yield studies &
Count of eligible studies, fractional and full &
OpenAlex, screened corpus &
Study &
2000--2026 &
Research-side outcome; 7,045 studies \\
\addlinespace
Study-location and crop coding &
Empirical study country; standardised crop &
Study text, coded &
Country-crop &
2000--2026 &
Allocates 2,031 studies to the panel \\
\addlinespace
Agricultural scale &
Production (t); harvested area (ha) &
FAOSTAT crop production &
Country-crop &
2018--2022 mean &
Denominator for research intensity and share comparisons \\
\addlinespace
Research need &
Nine indicators of food security, agricultural dependence, climate exposure, and production
instability; agricultural-scale shares recorded but excluded from the index &
FAOSTAT food security, macro, employment, and temperature domains &
Country &
Latest available &
Rank-aggregation index; four alternatives for sensitivity \\
\addlinespace
Scientific capacity &
R\&D expenditure; GDP per capita; tertiary enrolment; internet use; population &
World Bank indicators &
Country &
Latest available &
Model covariates; R\&D coverage governs the estimation sample \\
\addlinespace
Spatial data &
Country polygons and centroids &
Natural Earth admin-0 &
Country &
Static &
195 countries; basis for spatial weights \\
\bottomrule
\end{tabular}
\begin{tablenotes}
\item[] Reference periods refer to the observations entering the analysis rather than to the full extent of each source.
\end{tablenotes}
\end{table*}

\begin{table*}[!t]
\centering
\scriptsize
\setlength{\tabcolsep}{3.5pt}
\renewcommand{\arraystretch}{1.1}
\caption{Variable definitions, coverage, transformations, and missingness}
\label{tab:vardetail}
\begin{tabular}{@{}>{\raggedright\arraybackslash}p{4.6cm}
                  >{\raggedright\arraybackslash}p{3.3cm}
                  >{\raggedright\arraybackslash}p{1.95cm}
                  >{\raggedright\arraybackslash}p{1.55cm}
                  >{\raggedright\arraybackslash}p{1.95cm}
                  >{\raggedright\arraybackslash}p{3.15cm}@{}}
\toprule
Variable & Source & Period & Direction & Country coverage & Transformation and missingness \\
\midrule
\multicolumn{6}{@{}l}{\textit{Research-need components}} \\
Undernourishment (\%) & FAOSTAT food security & 2001--2024 & Higher & 172 & Percentile rank \\
Food insecurity, moderate or severe (\%) & FAOSTAT food security & 2015--2024 & Higher & 197 & Percentile rank \\
Cereal import dependency (\%) & FAOSTAT food security & 2001--2022 & Higher & 220 & Percentile rank \\
Dietary energy adequacy (\%) & FAOSTAT food security & 2001--2024 & Lower & 202 & Percentile rank \\
Employment in agriculture (\% of employment) & FAOSTAT employment & 1991--2025 & Higher & 226 & Percentile rank \\
Agricultural value added (\% of GDP) & FAOSTAT macro & 1990--2024 & Higher & 244 & Percentile rank \\
Agricultural value added per worker (USD) & FAOSTAT employment & 1991--2024 & Lower & 183 & Percentile rank \\
Temperature change since the 1951--1980 baseline ($^{\circ}$C) & FAOSTAT temperature & 1990--2025 & Higher & 285 in source, 190 in index & Percentile rank \\
Area-weighted yield volatility & Derived, FAOSTAT crops & 2018--2022 & Higher & 194 & Percentile rank \\
Share of global harvested area & Derived, FAOSTAT crops & 2018--2022 & --- & 199 & Recorded with the need data; excluded from the index (importance measure) \\
Share of global production & Derived, FAOSTAT crops & 2018--2022 & --- & 199 & Recorded with the need data; excluded from the index (importance measure) \\
\addlinespace
\multicolumn{6}{@{}l}{\textit{Scientific-capacity indicators, with missing panel cells}} \\
R\&D expenditure (\% of GDP) & World Bank & from 1996 & --- & 192 & Latest observation; 514 cells missing \\
GDP per capita, PPP & World Bank & from 1990 & --- & 242 & Logged; 154 cells missing \\
Tertiary enrolment (\%) & World Bank & from 1990 & --- & 241 & 125 cells missing \\
Internet users (\%) & World Bank & from 1990 & --- & 241 & 56 cells missing \\
Population, total & World Bank & from 1990 & --- & 260 & Logged; 56 cells missing; paging repair applied \\
Employment in agriculture (\%) & World Bank & from 1991 & --- & 230 & Retained; not in final specification \\
\addlinespace
\multicolumn{6}{@{}l}{\textit{Other panel quantities}} \\
Harvested area & FAOSTAT crops & 2018--2022 & --- & 199 & Logged; 112 cells missing \\
Yield volatility (CV) & Derived, FAOSTAT crops & 2018--2022 & --- & 199 & 248 cells missing; enters models only through the need index \\
\bottomrule
\end{tabular}
\begin{tablenotes}
\item[] Direction indicates whether a higher value corresponds to greater research need. Country coverage is the number of countries with at least one observation in the source extract. Missing cells are counts within the 2,616-cell panel for covariates entering the econometric models.
\end{tablenotes}
\end{table*}





\section{Methods and Procedures}

The analysis proceeds in four steps: eligible studies are allocated to country-crop systems to
measure research attention; attention is compared with agricultural scale and with
research need; the resulting country-level mismatch is tested for spatial dependence; and
regression models ask whether attention tracks need or scientific capacity once agricultural
size is held fixed. Table~\ref{tab:methods} maps each research question to the measure and
procedure used.

\subsection{Measuring Research Attention and Research Intensity}

The empirical unit is the country-crop cell. Let $i$ index countries and $k$ crops, and let
$s$ index eligible studies. A study examining $n^{C}_{s}$ countries and $n^{K}_{s}$ crops is
divided equally across the combinations it covers, so that research attention in cell $(i,k)$
is
\begin{equation}
R_{ik} \;=\; \sum_{s \in S_{ik}} w_{s},
\qquad
w_{s} \;=\; \frac{1}{n^{C}_{s}\, n^{K}_{s}},
\label{eq:frac}
\end{equation}
where $S_{ik}$ is the set of studies covering cell $(i,k)$. A cell with no eligible study takes $R_{ik}=0$; full counts are retained for sensitivity
analysis.

Research intensity normalises attention by the size of the agricultural activity at stake.
Writing $A_{ik}$ for mean harvested area and $P_{ik}$ for mean production over the reference
period,
\begin{equation}
I^{A}_{ik} \;=\; \frac{R_{ik}}{A_{ik}/10^{6}},
\qquad
I^{P}_{ik} \;=\; \frac{R_{ik}}{P_{ik}/10^{6}},
\label{eq:intensity}
\end{equation}
each defined only where the denominator is strictly positive, so higher values indicate more
research per unit of cultivated area or output. A complementary comparison places attention
and production on the same scale as shares of their global totals,
\begin{equation}
\Delta_{ik} \;=\; \frac{R_{ik}}{\sum_{jl} R_{jl}} \;-\; \frac{P_{ik}}{\sum_{jl} P_{jl}},
\label{eq:sharegap}
\end{equation}
so that $\Delta_{ik}>0$ marks a cell holding a larger share of the literature than of world
production. Both terms are shares of the global total across all cells, not shares within a
crop.

\subsection{Measuring Research Alignment and Mismatch}

We assess alignment against two distinct benchmarks. Equations~\eqref{eq:intensity} and
\eqref{eq:sharegap} compare attention with agricultural scale. The second comparison,
against research need, requires a country-level measure because the need index is defined for
countries rather than for country-crop systems. Cell-level attention is aggregated,
$R_{i}=\sum_{k} R_{ik}$, which preserves the fractional weights exactly.

Countries are then ranked on attention and on need, and the two ranks are differenced. Let
$\rho^{R}_{i}$ denote a country's percentile rank on $R_{i}$ and $\rho^{N}_{i}$ its percentile
rank on the research-need index, both expressed on a $0$--$100$ scale. The mismatch measure is
\begin{equation}
G_{i} \;=\; \rho^{R}_{i} \;-\; \rho^{N}_{i},
\label{eq:gap}
\end{equation}
measured in percentile points. Positive values mark countries better represented in the
literature than their need rank implies, negative values the reverse, and zero equal standing
in the two distributions. Percentile ranks place research attention and research need on a common scale. The measure
is defined only for countries with a computable need index and is missing rather than imputed
elsewhere. No composite combines importance, need, and capacity: each enters separately.

\subsection{Research-Need Measure}

The nine components are harmonised in direction so that a higher value always indicates
greater need, with dietary energy adequacy and agricultural value added per worker reversed
accordingly, then converted to percentile ranks across countries.

Rank aggregation is the primary index: the mean of a country's component percentile ranks,
re-expressed as a percentile on a $0$--$1$ scale among the 193 countries with at least five
observed components. Equal weighting, entropy weighting, the first principal component, and
the eleven-component variant are built for sensitivity analysis, though rank aggregation and
entropy weighting proved almost indistinguishable, so entropy weighting is not read as
independent corroboration. The component count is stored beside every value, and restricting
the index to countries with at least seven or nine observed components leaves the results in
Section~5 unchanged. Equation~\eqref{eq:gap} takes the index rescaled to $0$--$100$; the
models use the $0$--$1$ scale, so a coefficient reads as the change between the least and the
most needy country. Yield volatility enters research need through the index and appears
nowhere else in the models.

\subsection{Empirical Framework: Research Need versus Scientific Capacity}\label{sec:framework}

Two features of the outcome require separate models. Whether a country-crop system is studied
at all is a binary event, while how much it is studied is a non-negative count with a large
mass at zero. Participation and intensity are therefore estimated separately rather than
forced into one specification.

We model participation as
\begin{equation}
\Pr(y_{ik}=1 \mid \mathbf{x}_{ik}) \;=\; \Lambda\!\left(\beta_{0} + \mathbf{x}_{ik}'\boldsymbol{\beta}\right),
\qquad y_{ik} = \mathbf{1}\{R_{ik}>0\},
\label{eq:logit}
\end{equation}
where a probit link is estimated on the same sample
as an alternative. The covariate vector $\mathbf{x}_{ik}$ contains log harvested area, the
research-need index, research and development expenditure as a share of GDP, log GDP per
capita, tertiary enrolment, internet use, and log population. Log harvested area carries the
agricultural-importance channel and the need index the need channel; the remaining five are
capacity-related indicators rather than measures of a single capacity construct. Yield
volatility enters only through the need index, so no quantity appears on both sides of that
construction. Crop fixed effects are added in a robustness specification, so the
harvested-area association is also estimated from variation across countries within the same
crop.

We estimate intensity by Poisson pseudo-maximum likelihood,
\begin{equation}
\mathbb{E}\!\left[R_{ik} \mid \mathbf{x}_{ik}\right] \;=\;
\exp\!\left(\beta_{0} + \mathbf{x}_{ik}'\boldsymbol{\beta} + \ln A_{ik}\right),
\label{eq:ppml}
\end{equation}
where log harvested area enters as an offset with unit coefficient, so the model describes
studies per hectare and area leaves the covariate vector. Poisson pseudo-maximum likelihood accommodates the fractional counts \cite{SilvaTenreyro2006}. A
The unit elasticity that the offset imposes is tested rather than assumed: a specification
estimating the harvested-area coefficient freely is reported alongside, with a Wald test of
$\beta_{\text{area}} = 1$. A negative binomial model tests whether overdispersion argues against
the Poisson mean assumption, and a quasi-Poisson variance assumption is reported as a further
comparison.

Standard errors are clustered on country throughout, because cells within a country share
capacity covariates and are not independent draws. Dropping systems missing any covariate leaves 1,799 cells across 142 countries. Coefficients describe association: capacity, need,
and agricultural scale are jointly determined with research output, no instrument is used, and
no causal interpretation is offered.

\subsection{Spatial Analysis of Research Mismatch}

We evaluate spatial dependence on the country-level mismatch measure $G_{i}$ from
\eqref{eq:gap}, and on country aggregates of research output for comparison. A
row-standardised $k$-nearest-neighbour matrix with $k=6$ is primary because it leaves every
country with neighbours, whereas queen contiguity strands 43 islands that are not a random
subset of the world. Neighbours are found among polygon centroids projected to the Equal Earth
coordinate system. Because that projection preserves area rather than distance, the same
statistics are recomputed on true great-circle distance, which changes the neighbour set for 86
of 195 countries and leaves the results unchanged. Four alternatives are carried through the robustness suite, and the matrices are also
compared on a common connected subsample so that weight-matrix sensitivity is not confounded
with sample composition (\ref{app:robust}).

Global spatial dependence is measured by Moran's $I$ \cite{Moran1950},
\begin{equation}
I \;=\; \frac{n}{\sum_{i}\sum_{j} w_{ij}} \cdot
\frac{\sum_{i}\sum_{j} w_{ij}\,(x_{i}-\bar{x})(x_{j}-\bar{x})}{\sum_{i}(x_{i}-\bar{x})^{2}},
\label{eq:moran}
\end{equation}
Inference uses 9,999 conditional permutations.
Countries missing the variable under test are excluded and the weight matrix is rebuilt on
the observed sample, so no filled value enters a reported statistic.

Local Moran statistics \cite{Anselin1995} classify each country by whether a high or low
value sits among high or low neighbours, giving High--High, Low--Low, High--Low, and
Low--High categories. The $p$-values are
adjusted by the Benjamini--Hochberg procedure \cite{BenjaminiHochberg1995} at $\alpha=0.05$
and a country is called a cluster only on the adjusted value. Getis--Ord $G^{*}_{i}$
\cite{GetisOrd1992} is computed alongside under the same scheme. Local statistics use the
primary matrix only. Residual spatial dependence is tested by applying \eqref{eq:moran} to
country-mean residuals from both models.

\subsection{Robustness and Sensitivity Analysis}

Fifteen variants test whether measured concentration depends on counting rules, document type,
location-coding confidence, the allocation scope, the method definition, the largest contributing
countries, or the period examined. Fourteen re-estimate the participation model across country
exclusions, income groups, crop and region fixed effects, and alternative need indices. Spatial
results are re-computed across neighbour definitions, a common connected subsample, and both need
indices. Results are reported rather than graded. For each quantity Table~\ref{tab:robust} gives the
baseline estimate, the range across variants, and how many variants retain the sign and exclude
zero. Thresholds separating stable from unstable findings were not fixed before the results were
examined, so no categorical label is attached to them and the spread is given instead. Quantities
whose intervals contain zero are reported as such rather than omitted.
\ref{app:robust} lists the specifications.

\begin{table*}[!t]
\centering
\footnotesize
\setlength{\tabcolsep}{4pt}
\renewcommand{\arraystretch}{1.1}
\caption{Empirical measures and analytical procedures}
\label{tab:methods}
\begin{tabular}{@{}>{\raggedright\arraybackslash}p{3.6cm}
                  >{\raggedright\arraybackslash}p{4.3cm}
                  >{\raggedright\arraybackslash}p{1.9cm}
                  >{\raggedright\arraybackslash}p{4.7cm}@{}}
\toprule
Question & Measure & Unit & Procedure \\
\midrule
Where is research done? & Fractional study count, Eq.~\eqref{eq:frac} & Country-crop &
Allocation and description \\
Is attention proportional to agriculture? & Intensity and share difference,
Eqs.~\eqref{eq:intensity}--\eqref{eq:sharegap} & Country-crop &
Normalisation against harvested area and production \\
Is attention aligned with need? & Mismatch measure, Eq.~\eqref{eq:gap} & Country &
Percentile ranking and differencing \\
Need or capacity? & Participation and intensity,
Eqs.~\eqref{eq:logit}--\eqref{eq:ppml} & Country-crop &
Logit and probit; Poisson pseudo-maximum likelihood with area offset; clustered by country \\
Is mismatch clustered? & Moran's $I$ and local statistics, Eq.~\eqref{eq:moran} & Country &
Permutation inference with Benjamini--Hochberg adjustment \\
\bottomrule
\end{tabular}
\end{table*}

\begin{table*}[!t]
\centering
\caption{Summary statistics for the variables entering the regression models}
\label{tab:summstats}
\begin{tabular*}{\textwidth}{@{\extracolsep\fill}l r r r r r r r r@{}}
\toprule
Variable & N & Mean & SD & Median & Min & Max & Mean, $y=1$ & Mean, $y=0$ \\
\midrule
Log harvested area & 1799 & 10.215 & 3.048 & 10.499 & 0.000 & 17.621 & 12.884 & 9.333 \\
R\&D expenditure (\% GDP) & 1799 & 0.806 & 1.051 & 0.392 & 0.006 & 6.346 & 1.215 & 0.671 \\
Research-need index (0-1) & 1799 & 0.513 & 0.283 & 0.523 & 0.016 & 1.000 & 0.399 & 0.551 \\
Log GDP per capita & 1799 & 9.599 & 1.087 & 9.726 & 6.972 & 11.786 & 9.941 & 9.486 \\
Tertiary enrolment (share) & 1799 & 0.467 & 0.326 & 0.455 & 0.019 & 1.651 & 0.579 & 0.430 \\
Internet use (share) & 1799 & 0.705 & 0.261 & 0.820 & 0.086 & 1.000 & 0.775 & 0.682 \\
Log population & 1799 & 16.835 & 1.632 & 16.903 & 11.512 & 21.104 & 17.721 & 16.543 \\
Fractional study count & 1799 & 1.084 & 6.641 & 0.000 & 0.000 & 159.519 & 4.361 & 0.000 \\
\bottomrule
\end{tabular*}
\begin{tablenotes}
\item[] The estimation sample is 1,799 country-crop cells across 142 countries. $y=1$ marks cells carrying at least one eligible study (447 cells); $y=0$ marks cells carrying none (1,352 cells).
\end{tablenotes}
\end{table*}

\section{Results}

\subsection{Global Distribution of AI/ML Crop-Yield Research}

Research attention concentrates in a small number of countries and crops. Of the 199
countries in the panel, 122 carry at least one eligible study and 77 carry none, and only 485
of the 2,616 country-crop systems contain any research at all (Figure~\ref{fig:map}). The Gini
coefficient of fractional study counts across all 199 countries is 0.848, the principal
concentration measure because the 77 countries with no resolved study are part of the
distribution being described. Conditioning on the 122 countries that carry a study gives 0.750, reported for comparison.

The distribution is top-heavy at every level. The most-studied country accounts for 16.1\% of
fractional output, the top three for 43.0\%, and the top twenty for 73.4\%. The United States, India, and China lead with 327.7, 292.7, and 252.3 fractional
studies. By income group, upper-middle-income countries hold 34.1\% of fractional output and
high-income countries 33.6\%, against 4.2\% for low-income countries.

Crop coverage is similarly narrow. Wheat, maize, and rice together take 69.3\% of fractional crop attention (566.4, 451.2, and
390.5 fractional studies), and the Gini across crops is 0.719 (Figure~\ref{fig:crops}, \ref{app:figs}). The literature is also recent: the median
publication year is 2023 and 80.8\% of eligible studies appeared in 2018 or later.

\subsection{Research Attention Relative to Agricultural Scale}

Research attention tracks agricultural scale only loosely at the country-crop level.
Harvested area is the cross-crop benchmark, because a hectare is comparable across crops
while a tonne of sugarcane and a tonne of wheat are not; production is compared only within
a crop. The Spearman correlation between fractional study counts and harvested area is 0.511
across cells, rising to 0.756 when both are aggregated to countries (0.506 and 0.791 for
production), so the correspondence that appears at country level weakens once the crop
dimension is restored (Figure~\ref{fig:share}).

Departures from proportionality run in both directions. Against harvested area, the three
major cereals absorb more attention than the land they occupy: wheat holds 28.0\% of
fractional attention on 17.9\% of panel area, maize 22.3\% on 16.5\%, and rice 19.3\% on
13.6\%. Oil palm is the clearest shortfall: 0.4\% of attention on 2.2\% of area, and
Indonesia, with 48.5\% of world oil-palm area, carries no eligible oil-palm study at all.
Within single crops, production-based comparisons point the same way: Brazil holds 24.5\%
of world sugarcane research against 38.7\% of world sugarcane production, and China 12.2\%
of rice research against 27.3\% of rice production, while United States maize and Chinese
wheat run above their within-crop production shares. Comparisons that pool production
tonnage across crops overstate the shortfall for high-biomass crops; sugarcane holds a
quarter of panel tonnage but 2.2\% of panel area, and against area its attention share of
3.9\% is no longer a deficit.

\subsection{Research Need and Scientific Capacity}

Research attention differs from the geographic ranking of research need.
Ranking countries on fractional study counts and on the rank-aggregation need index and
differencing the two percentiles gives a mismatch measure defined for 193 countries, with 56
falling in the high-need, low-research quadrant (Figure~\ref{fig:need}); the eleven-component
variant that adds the agricultural-scale shares back gives 48, because large producers then
count as needier. The descriptive mismatch reported here is distinct from the conditional association
estimated in Section~\ref{sec:reg}. The countries with the largest shortfalls are concentrated
among low-income economies, and 71 of the 193 tie at the lowest research percentile because
none carries an eligible study.

Attention is instead concentrated among higher-capacity countries. The 4.2\% of fractional
output held by low-income countries sits against their far larger share of the countries in
the panel, and the same countries occupy the lower end of the research and development
expenditure distribution. Whether that association survives controls for agricultural scale is
the question the regressions address.

\subsection{Regression Results}\label{sec:reg}

Agricultural scale and several capacity indicators are associated with where research is
conducted, whereas the association with measured research need is imprecisely estimated.
Table~\ref{tab:reg} reports the participation logit and the intensity model, estimated on
1,799 and 1,792 country-crop systems across 142 countries, with standard errors clustered by
country.

Harvested area has the largest and most precisely estimated association with whether a
country-crop system is studied at all, with a logit coefficient of 0.713 (95\% CI 0.613 to
0.813). The comparison across covariates does not rest on raw coefficients: an
interquartile-range increase in log harvested area raises the predicted probability of
participation by 35 percentage points, about five times the corresponding change for any
other covariate, the largest of which is 7
points for log GDP per capita. Average marginal effects agree: harvested area at 0.081 ($p < 0.001$) and research and
development expenditure at 0.036 ($p = 0.010$) are the only covariates distinguishable from zero
(Table~\ref{tab:ame}). Figure~\ref{fig:predprob} plots the fitted participation probability
across the observed range of harvested area. Adding crop fixed effects, so that the association is identified from variation across
countries within the same crop, moves the coefficient to 0.656, and estimating the model
separately for each of the five most studied crops returns coefficients from 0.741 to 1.168, all
with intervals excluding zero (Table~\ref{tab:specchecks}). Research and development expenditure is also
associated with participation, at 0.315 (CI 0.074 to 0.556), and with intensity once
harvested area is held fixed as an exposure offset, at 0.522 (CI 0.281 to 0.762); an
interquartile-range increase multiplies expected intensity by 1.5. Internet use and tertiary
enrolment are distinguishable from zero only in the intensity model, internet use positively
(2.179, CI 0.543 to 3.648) and tertiary enrolment negatively ($-0.985$, CI $-1.817$ to
$-0.153$). Log GDP per capita and log population are not distinguishable from zero in either
model.

The need association is not distinguishable from zero once the intensity model drops a
restriction the data reject. In the participation model the need coefficient is 0.944 with a
standard error of 0.961 and a 95\% interval from $-0.940$ to 2.829: positive but imprecisely
estimated, an interval wide enough to contain a large association in either direction. In the
intensity model with harvested area entered as an exposure offset it is 1.261 with an interval
from 0.105 to 2.417, and it stays marginally distinguishable from zero with crop fixed effects
(1.221, CI 0.025 to 2.416). That specification fixes the harvested-area elasticity at one.
Estimated freely, the elasticity is 0.849 (SE 0.049, CI 0.754 to 0.944), and a Wald test rejects
unit elasticity ($\chi^{2}_{1} = 9.70$, $p = 0.002$); under the free specification the need
coefficient falls to 0.950 with an interval from $-0.186$ to 2.086. No specification therefore
distinguishes the need association from zero, which is a statement about the precision these
data support and not evidence that no association exists. Index construction affects the size of the estimate without changing that conclusion: under the
eleven-component variant, which contains the agricultural-scale shares, the two need coefficients
are $-0.016$ and 0.209, against 0.944 and 1.261 under the scale-excluded measure.
Figure~\ref{fig:coef} in \ref{app:figs} plots every estimate with its interval.

The participation model achieves a McFadden pseudo-$R^{2}$ of 0.366 on 447 events. The
intensity model retains 1,792 of the 1,799 cells because the exposure offset $\ln A_{ik}$ is
undefined for seven covariate-complete cells recording zero harvested area. The counts are overdispersed, with a conditional variance about forty times the mean
(overdispersion ratio $\approx 40$). Alternative variance assumptions widen the intervals without
changing the point estimates, and a negative binomial returned no standard errors under the
clustered covariance estimator; the dispersion diagnostics are reported in
Table~\ref{tab:diagnostics}. Neither model shows significant residual spatial autocorrelation. Moran's $I$ on country-mean
residuals is $-0.029$ ($p = 0.324$) for participation and $-0.020$ ($p = 0.386$) for intensity,
and Lagrange Multiplier tests reject neither a spatial lag nor a spatial error process for either
model under the primary weights (Table~\ref{tab:diagnostics}).

\subsection{Spatial Clustering of Research Mismatch}

The mismatch between research standing and need is spatially clustered. Global Moran's $I$
on the country-level mismatch measure, computed over the 193 countries with an observed
value, is 0.363 ($p = 0.0001$, 9,999 permutations) under the primary six-nearest-neighbour
matrix. The statistic moves little with the neighbour definition, from 0.358 to 0.378 across the
$k$-nearest-neighbour alternatives and 0.362 under queen contiguity on its connected subsample,
and holding the sample fixed at that connected subsample gives 0.400 to 0.421 across matrices,
so matrix sensitivity is small once sample composition is held constant. Rebuilding the
neighbour sets on great-circle rather than projected distance changes which countries are
neighbours for 86 of 195 cases and leaves the statistic at 0.364. The measure also survives the
tie convention: 71 of the 193 countries tie at the lowest research percentile, and the statistic
runs from 0.291 to 0.371 across four rank conventions and reaches 0.187 under a standardised
log-count alternative that preserves zeros, with $p = 0.0001$ throughout.
Under the eleven-component index the same observed-data statistic is 0.394. Mismatch is the
most strongly clustered of the quantities examined: the log fractional study count returns
0.157 ($p = 0.0007$) and crop-system coverage 0.143 ($p = 0.0006$), while the
citation-weighted count returns 0.022 ($p = 0.088$) and is not significant under the primary
matrix.

Local statistics identify clustering in both directions. After Benjamini--Hochberg
adjustment, 27 countries carry a significant local statistic: 17 form High--High clusters, 6
Low--Low, 2 High--Low, and 2 Low--High (Figure~\ref{fig:lisa}). All 17 High--High countries
lie in western and northern Europe. The Low--Low cluster holds the Republic of the
Congo, the Democratic Republic of the Congo, Gabon, Sudan, Djibouti, and Kiribati; Ethiopia
is a High--Low outlier, better represented than its low-mismatch neighbours. Getis--Ord
$G^{*}_{i}$ agrees in direction, returning 19 hot spots and 8 cold spots. The
eleven-component index gives the same geography with a broader African Low--Low set (19
High--High, 11 Low--Low).

Research output on its own clusters differently. Its local statistics after adjustment give 21
Low--Low clusters against 2 High--High, and Getis--Ord returns 3 hot spots against 21 cold
spots, so raw output clusters as shared absence while the need-adjusted mismatch clusters in
both directions.

\subsection{Robustness and Sensitivity Results}

Table~\ref{tab:robust} gives each estimated quantity, its baseline, and its spread across the
robustness suite. Concentration is stable across publication periods (Figure~\ref{fig:temporalgini}). Splitting
the outcome by publication period leaves the area association unchanged, at 0.750 for studies
published before 2020 and 0.748 for those published from 2020 onward, while the research and development coefficient falls from 0.567 to 0.285 and the need
coefficient stays imprecise in both periods (Table~\ref{tab:periods}).

\begin{table*}[!t]
\centering
\footnotesize
\setlength{\tabcolsep}{4pt}
\renewcommand{\arraystretch}{1.12}
\caption{Estimated quantities across the robustness suite}
\label{tab:robust}
\begin{tabular}{@{}>{\raggedright\arraybackslash}p{5.1cm}
                  r c r >{\raggedright\arraybackslash}p{3.9cm}@{}}
\toprule
Estimated quantity & Baseline & Range across variants & Variants & Sign retained / interval excludes zero \\
\midrule
Gini of country research counts, all countries & 0.847 & [0.772, 0.899] & 16 & 16 of 16 / not applicable \\
Logit coefficient on log harvested area & 0.713 & [0.654, 0.861] & 14 & 14 of 14 / 14 of 14 \\
Global Moran's $I$ on the mismatch measure & 0.363 & [0.187, 0.425] & 14 & 14 of 14 / 14 of 14 \\
Share of crop attention held by wheat, maize and rice & 69.3\% & 67.7\% under the machine-learning subset & 2 & not applicable \\
Logit coefficient on R\&D expenditure & 0.315 & [$-0.074$, 2.218] & 14 & 12 of 14 / 8 of 14 \\
Global Moran's $I$ on citation-weighted output & 0.022 & [0.014, 0.147] & 5 & 5 of 5 / 1 of 5 \\
Logit coefficient on the research-need index & 0.944 & [$-0.348$, 0.950] & 8 & 6 of 8 / 0 of 8 \\
PPML coefficient on the research-need index & 1.261 & [$-0.118$, 1.568] & 6 & 5 of 6 / 2 of 6 \\
Harvested-area elasticity, freely estimated & 0.849 & [0.754, 0.944] & 1 & unit elasticity rejected, $p=0.002$ \\
\bottomrule
\end{tabular}
\begin{tablenotes}
\item[] Baseline is the value in the main specification: the nine-component scale-excluded need
index, observed-sample spatial statistics, and harvested area entered as an exposure offset in the
intensity model. Range is the span across the variants listed in \ref{app:robust}, which include
crop and region fixed effects, both need indices, the corrected coverage floor, allocation-scope
and method-definition restrictions, alternative weight matrices, great-circle neighbours,
alternative tie conventions, and the free-elasticity intensity specification. The final column
reports how many variants preserve the sign of the baseline estimate and how many return an
interval excluding zero; no categorical stability label is assigned, because the thresholds that
would define one were not fixed before the results were examined. Concentration variants are
computed on the study-level allocation, whose fractional total is 2,031 against 2,020.2 inside the
2,616-cell panel, so the baseline Gini differs from the panel figure of 0.848 in the third decimal.
\end{tablenotes}
\end{table*}

Concentration is the most durable result. Across the 14 variants that recompute it, the
all-country Gini stays between 0.772 and 0.899 against a baseline of 0.847, and removing the
United States, China, and India together still leaves 0.772. Journal-only estimation returns
0.849 on 1,712 studies, so conference proceedings do not drive the pattern, and excluding
cells whose location rests on the weaker locative cue returns 0.849. Restricting allocation
to studies whose overall scope names specific countries gives 0.839, and restricting the
corpus to studies whose title or abstract carries an explicit machine-learning method term
gives 0.870, so neither the allocation rule nor the breadth of the method definition drives
the concentration. The conditional Gini over countries that carry research moves between
0.619 and 0.836 across the same variants.

The area result holds as widely: across 14 model variants, including crop and region fixed
effects and both need indices, the logit coefficient on log harvested area stays between 0.654
and 0.861 and is positive with an interval excluding zero in every one. The research and
development coefficient is less settled, ranging from $-0.074$ to 2.218 and moving materially
across income-group splits, so it excludes zero in 8 of the 14 variants. Clustering of the
mismatch measure holds across neighbour definitions, index variants, great-circle distance, and
tie conventions. Clustering of the citation-weighted count ranges from 0.014 to 0.147 and
excludes zero in only 1 of 5 variants, so no clustering claim is made for it. The need
association with participation has a standard error near one in every variant and excludes zero
in none; the need association with intensity is positive under the primary index with and without crop
fixed effects, but indistinguishable from zero under the eleven-component variant and under the
free-elasticity specification. Adding the capacity indicators to a model already containing agricultural scale and research
need improves fit, whereas adding research need to scale and capacity does not
(Table~\ref{tab:blocks}). Aggregating to countries reproduces the pattern only in part
(Table~\ref{tab:countrylevel}). A
country-level Poisson model with an area offset keeps the research and development association
positive (0.278, CI 0.069 to 0.487) and leaves need imprecise (0.533, CI $-0.638$ to 1.704). A
log-count specification instead returns a negative need coefficient (\(-1.083\), CI $-2.072$ to
$-0.094$), an imprecise research and development coefficient, and log population as the dominant
covariate.

\begin{table*}[!t]
\centering
\footnotesize
\setlength{\tabcolsep}{5pt}
\renewcommand{\arraystretch}{1.12}
\caption{Determinants of research participation and research intensity}
\label{tab:reg}
\begin{tabular*}{\textwidth}{@{\extracolsep\fill}l r@{\hspace{6pt}}l r@{\hspace{6pt}}l@{}}
\toprule
& \multicolumn{2}{c}{Participation (logit)} & \multicolumn{2}{c}{Intensity (PPML)} \\
\cmidrule(lr){2-3}\cmidrule(lr){4-5}
Covariate & Estimate & 95\% CI & Estimate & 95\% CI \\
\midrule
Log harvested area & 0.713 & [0.613, 0.813] & \multicolumn{2}{c}{offset} \\
 & (0.051) & & & \\
R\&D expenditure (\% GDP) & 0.315 & [0.074, 0.556] & 0.522 & [0.281, 0.762] \\
 & (0.123) & & (0.123) & \\
Research-need index & 0.944 & [$-0.940$, 2.829] & 1.261 & [0.105, 2.417] \\
 & (0.961) & & (0.590) & \\
Log GDP per capita & 0.347 & [$-0.331$, 1.026] & $-0.052$ & [$-0.519$, 0.415] \\
 & (0.346) & & (0.238) & \\
Tertiary enrolment & 0.257 & [$-0.663$, 1.176] & $-0.985$ & [$-1.817$, $-0.153$] \\
 & (0.469) & & (0.425) & \\
Internet use & 1.312 & [$-0.665$, 3.290] & 2.179 & [0.543, 3.815] \\
 & (1.009) & & (0.835) & \\
Log population & 0.111 & [$-0.100$, 0.321] & $-0.108$ & [$-0.218$, 0.002] \\
 & (0.107) & & (0.056) & \\
\midrule
Country-crop cells & \multicolumn{2}{c}{1,799} & \multicolumn{2}{c}{1,792} \\
Countries & \multicolumn{2}{c}{142} & \multicolumn{2}{c}{142} \\
Pseudo-$R^{2}$ & \multicolumn{2}{c}{0.366} & \multicolumn{2}{c}{---} \\
Residual Moran's $I$ & \multicolumn{2}{c}{$-0.029$ ($p=0.324$)} & \multicolumn{2}{c}{$-0.020$ ($p=0.386$)} \\
\bottomrule
\end{tabular*}
\begin{tablenotes}
\item[] The participation column reports a logit for whether a country-crop cell carries any eligible study. The intensity column reports Poisson pseudo-maximum-likelihood estimates of the fractional study count with log harvested area as an exposure offset. The research-need index is the nine-component scale-excluded measure; yield volatility enters through the index and is not a separate regressor. Standard errors, clustered by country, are in parentheses; 95\% confidence intervals are in brackets.
\end{tablenotes}
\end{table*}

\begin{table*}[!t]
\centering
\caption{Average marginal effects on the probability that a country-crop system is studied}
\label{tab:ame}
\begin{tabular*}{\textwidth}{@{\extracolsep\fill}l r r r r r r@{}}
\toprule
Covariate & AME & (SE) & 95\% CI & 25th pct. & 75th pct. & IQR effect (pp) \\
\midrule
Log harvested area & 0.081 & (0.004) & [0.072, 0.089] & 8.27 & 12.34 & 35.0 \\
R\&D expenditure (\% GDP) & 0.036 & (0.014) & [0.008, 0.063] & 0.18 & 0.94 & 2.7 \\
Research-need index (0-1) & 0.107 & (0.109) & [$-$0.108, 0.321] & 0.26 & 0.76 & 5.2 \\
Log GDP per capita & 0.039 & (0.040) & [$-$0.038, 0.117] & 8.81 & 10.53 & 7.0 \\
Tertiary enrolment (share) & 0.029 & (0.053) & [$-$0.075, 0.133] & 0.15 & 0.71 & 1.6 \\
Internet use (share) & 0.149 & (0.111) & [$-$0.070, 0.367] & 0.55 & 0.91 & 5.6 \\
Log population & 0.013 & (0.012) & [$-$0.011, 0.036] & 15.76 & 17.82 & 2.6 \\
\bottomrule
\end{tabular*}
\begin{tablenotes}
\item[] Average marginal effects from the participation logit of Table~\ref{tab:reg}, computed by the delta method with standard errors clustered by country. The final column reports the change in mean predicted probability, in percentage points, when the covariate moves from its 25th to its 75th percentile with all other covariates held at their observed values. Percentage-point changes are comparable across covariates; raw coefficients are not.
\end{tablenotes}
\end{table*}





\begin{figure*}[!t]
\centering
\includegraphics[width=0.98\linewidth,height=0.68\textheight,keepaspectratio]{fig_01_research_intensity_map.pdf}
\caption{Fractional AI/ML crop-yield research output by country. The map shows the 195 countries
in the boundary layer; four of the 199 panel countries have no matching geometry and none of the
four carries an eligible study, so 73 zero-study countries appear here against 77 in the full
panel. Countries with no resolved eligible study are shown in the zero class and are
distinguished from countries with no available data.}
\label{fig:map}
\end{figure*}

\begin{figure*}[!t]
\centering
\includegraphics[width=0.62\linewidth,height=0.80\textheight,keepaspectratio]{fig_A6_sample_flow.pdf}
\caption{Study identification and analytical sample flow. Counts on the right are the records
surviving each stage; notes on the left of each arrow give what the stage removes. The
fractional totals differ because 10.9 study-equivalents fall on country-crop pairs outside the
2,616-cell panel.}
\label{fig:flow}
\end{figure*}

\begin{figure*}[!t]
\centering
\includegraphics[width=0.98\linewidth,height=0.68\textheight,keepaspectratio]{fig_03_v2_scale_alignment.pdf}
\caption{Share of fractional research against share of world harvested area, for
country-crop systems (left) and countries (right). Both axes are logarithmic; the dashed
line marks equal shares; units with no resolved study carry no position on a logarithmic axis and are drawn as hollow
triangles on the shaded band below the plotting range, which is off scale rather than a small
positive value.}
\label{fig:share}
\end{figure*}

\begin{figure*}[!t]
\centering
\includegraphics[width=0.98\linewidth,height=0.68\textheight,keepaspectratio]{fig_04_v2_need_vs_research.pdf}
\caption{Country research percentile against research-need percentile. 193 countries with
a computable nine-component need index. The shaded quadrant holds the 56 countries at or
above the median for need and at or below the median for research.}
\label{fig:need}
\end{figure*}

\begin{figure*}[!t]
\centering
\includegraphics[width=0.98\linewidth,height=0.68\textheight,keepaspectratio]{fig_05_v2_mismatch_lisa.pdf}
\caption{Local Moran clusters of the research-need mismatch measure. Categories are assigned on
Benjamini--Hochberg adjusted pseudo $p$-values from 9,999 permutations, computed over the 193
countries with an observed mismatch value under the primary six-nearest-neighbour weights. The
dashed rectangles mark the extents enlarged in the two panels below, which label every clustered
country. Panel (b) shows five of the six Low--Low countries; the sixth, Kiribati, lies in the
central Pacific.}
\label{fig:lisa}
\end{figure*}

\begin{figure*}[!t]
\centering
\includegraphics[width=0.86\linewidth,height=0.52\textheight,keepaspectratio]{fig_A3_predicted_probability.pdf}
\caption{Predicted probability that a country-crop system carries research, as a function of log
harvested area, with all other covariates held at their sample means. The band is a 95\%
confidence interval; the histogram shows the distribution of log harvested area across the
estimation sample.}
\label{fig:predprob}
\end{figure*}

% >>> DISCUSSION-BEGIN
\section{Discussion}

\subsection{Main Findings and Their Interpretation}

AI-based crop-yield research aligns in part with agricultural scale and unevenly with
selected capacity indicators; whether it aligns with research need cannot be resolved at the
precision available.
Harvested area is the one covariate that behaves the same way in every specification: the more
land a country devotes to a crop, the more likely that country-crop system is to have been
studied. The coefficient stays within a narrow band across the model variants, including crop
and region fixed effects, and it is positive with an interval excluding zero in each of the five
most studied crops estimated separately. Cropland
extent accounts for a real share of where evidence exists, and a reading of the research map
that ignored it would mistake agricultural structure for neglect.

Alignment stops short of proportionality. The rank correlation between fractional study counts
and agricultural scale weakens once the crop dimension is restored, so agreement
visible in national totals dissolves at the level where crops are grown and modelled, and
research share rises with area at roughly a third of the rate proportionality would require
(Section~5.2). The departures run in both directions: the major cereals absorb more attention
than the land they occupy, whereas oil palm is the clearest shortfall (Figure~\ref{fig:crops}).
Neither pattern is visible in national totals alone.

The association with research need is estimated much less precisely. The intensity association
is positive under the scale-excluded index, but the result is sensitive to both index
construction and the harvested-area restriction: it disappears under the eleven-component
variant, and estimating the elasticity freely returns an interval containing zero. Imprecision of
this kind does not show that research need is unrelated to research attention. The need
association remains unidentified, and that is the limit of what these data support.

Capacity indicators are estimated more precisely than need but show no consistent directional
pattern. R\&D
expenditure and internet access are positively associated with research intensity, whereas
tertiary enrolment enters negatively and GDP per capita is imprecisely estimated. The mixed
signs caution against interpreting the capacity covariates as a single latent dimension: what
the estimates support is an association with particular capacity-related indicators,
principally research financing, not with capacity in general. The expenditure association also moves across income-group splits, so it carries less weight
than the area result.

The mismatch between a country's research standing and its need rank is spatially clustered
more strongly than any other quantity examined, and it clusters in both directions where
research output alone does not. Output on its own produces mostly Low--Low clusters, a pattern dominated by neighbouring
countries with low research output; mismatch produces High--High clusters across western and
northern Europe against Low--Low clusters concentrated in sub-Saharan Africa
(Figure~\ref{fig:lisa}).
Over- and under-representation relative to need are regional conditions shared with
neighbours rather than isolated national cases.

\subsection{Geographic Representation and Transferability}

Concentration matters if models do not transfer across settings. The transferability studies reviewed above answer in a consistent direction:
apparent accuracy degrades as the evaluation domain moves away from the training domain, and
the degradation is invisible under the validation designs most commonly reported
\cite{Habibi2024,Okupska2026,Stiller2024}.

Those comparisons are made within fields and regions, and the present study does not estimate
model transfer across countries. Spatially concentrated evidence does not by itself establish
that models perform badly in settings the literature has not covered. The concern is narrower
and measurable: of the 2,343 country-crop systems with a positive harvested-area denominator,
1,864 carry no eligible study, and those systems account for 19.1\% of the panel's harvested
area. Almost four-fifths of the world's cropping systems, covering close to a fifth of its
cropland, have no in-domain evidence to degrade from. The concentration limits the geographic breadth of the evidence base; it does
not demonstrate failure of any model, and this design could not detect such failure.

\subsection{Research Need versus Scientific Capacity}

Research need and scientific capacity answer different questions, and the analysis keeps them
apart: need describes where additional evidence would be useful, capacity where research
can be produced. Concentration among a small group of countries is not specific to
this literature \cite{Bueno2025}; what the country-crop unit adds is the ability to ask
which of the two it follows.

The observed geography is more consistent with differences in some dimensions of scientific
capacity, particularly research financing, than with differences in need. Support for that reading is asymmetric: the R\&D association is estimated with a standard error
of about a tenth of a log-odds but moves across income-group splits, while the need association
is imprecise for participation and, for intensity, positive only under a restriction the data
reject. The available precision does not support a stronger conclusion.

The asymmetry matches what the agricultural research literature reports about the systems
producing the work. In low- and middle-income countries the smallest agricultural research systems sit inside
weaker surrounding institutions and thinner innovation environments \cite{NinPratt2023},
a constraint on producing research that operates independently of how much a country would
gain from the evidence.
\citet{NinPratt2021} adds a caution: the conventional intensity ratio understates the
research effort developing countries make, so low measured intensity should not be read as low
effort, even though an investment gap close to half of R\&D spending remains. Low output in the Low--Low cluster is therefore better read as a
constraint on research production than as a revealed judgement about the value of the
research, particularly because \citet{Rosegrant2023} estimate a central benefit--cost
ratio of 33 for expanded agricultural R\&D in the Global South.

Food-insecure countries receive about a fifth of global agricultural investment
\cite{Zhao2023}, a distribution that resembles the research pattern reported here. The
resemblance is untested, because the analysis does not link investment flows to study locations. The Low--Low cluster should not be read as uniform
either: returns to agricultural innovation investment vary by subregion within sub-Saharan
Africa \cite{Malec2024}.

Agricultural scale on its own does not account for research geography either. Harvested
area is associated with whether a system is studied but not with which large systems attract
attention, because the departures from proportionality noted above occur among the largest
producers.

\subsection{Implications for AI-Based Crop-Yield Research}

Study location deserves to be reported as a matter of course. Author affiliation does not
identify where a model was trained or evaluated, and fewer than half of the eligible studies
here carry a usable empirical location (\ref{app:coding}). An explicit statement of study
area would make the geography of this literature measurable without inference.

Geographically independent validation is what would establish whether models generalise across
geographic domains. The
transferability studies above measure that gap within regions; extending the practice to
settings where evidence is thin would turn a question about generalisability into a measured
quantity.

Low research intensity relative to production or need marks where evidence is absent, not
where funding should go. The mismatch measure ranks evidence available against evidence that
would be useful, and says nothing about the cost of producing it, the capacity to absorb it,
or competing claims on research budgets. Reading the Low--Low cluster as an allocation rule
would extend the measure past what it estimates.

Comparing counts against agricultural denominators is more informative than publication
totals. Raw output confounds the size of a country's research system with the size of its
agriculture. Agricultural denominators distinguish research concentration associated with
agricultural scale from disproportionate allocation.


\subsection{Limitations}

Location resolution bounds every geographic statement made here. Fewer than half of the
eligible studies (43.4\%) carry a usable empirical location, and about a quarter resolve to
specific countries (\ref{app:coding}). The geographic patterns therefore describe the resolvable
evidence base rather than every eligible publication. Coding accuracy was measured rather than
assumed: allocation to countries is correct for 90.4\% of an adjudicated probability sample
(95\% CI 84.3 to 94.3), but 62.0\% of sampled unresolved records name a study area the coder did
not recover, which projects to roughly a thousand studies. Resolved and unresolved studies differ on every observable tested (Table~\ref{tab:resolution}):
resolved studies carry more citations (23 against 15), are more likely to have an abstract
(93.9\% against 80.2\%), and are less often conference items (15.7\% against 23.4\%), while
publication years differ by half a year. Resolved studies are also somewhat less likely to name
an explicit machine-learning method (49.8\% against 62.4\%), so the resolvable base is not a
methodologically narrower slice of the corpus. Against that, the countries named
in the unrecovered records track the existing allocation closely by income group (high income
34.5\% against 33.6\%, low income 3.4\% against 3.6\%) and by mean need percentile (0.380 against
0.355), so on the dimensions that can be observed the under-resolution does not look
geographically selective. It may still be selective on dimensions that cannot.

Screening was automated, with a false-inclusion rate of 10.0\% and a false-exclusion rate of
2.2\% in the largest exclusion stratum, both estimated from 85 adjudicated records whose interval upper bounds reach 0.236 for false
inclusion and 0.116 for false exclusion (\ref{app:corpus}). Counts carry
classification error that is estimated but not tightly bounded, and a cell recorded as zero could
reflect a missed or unresolved study rather than an absent literature. Both the screening and the
coding validations were adjudicated by model-assisted review against the same criteria rather
than by independent human reading, so they measure the extractor against its own evidence base
and require human confirmation before they can be treated as settled.

Research need is measured at country level while the analytical unit is the country-crop
system, so the need index cannot vary across crops within a country. The need coefficient is
identified from between-country variation alone, which is part of why it is imprecise, and a
need measure varying by crop might behave differently. The covariates are the most recent observations rather than a common year. The models therefore
compare the accumulated 2000--2026 evidence stock with current conditions, and are read as
measures of present alignment rather than as historical determinants of research production. The
median reference year is 2023 or later for every covariate, although R\&D expenditure is older in
low- and lower-middle-income countries, with medians of 2018 and 2019 against 2023 for the two
higher-income groups (Table~\ref{tab:diagnostics}). Missing capacity data compounds the restriction: R\&D expenditure
is the sparsest covariate (Table~\ref{tab:vardetail}), and the models are estimated where
capacity indicators are observed. Observation is itself capacity-related, so the countries
least represented in the estimation sample are disproportionately those the need question
concerns most.

Study counts measure research attention rather than research quality or societal value, and
publication output is at best a crude indication of underlying research capacity
\cite{CashGibson2018}. No specification identifies why a system is studied, and reverse causation is not excluded. Spatial results depend
on the neighbour structure imposed on a country set containing many islands, though the
mismatch clustering holds across the neighbour definitions, the common connected subsample,
and both need indices.
% <<< DISCUSSION-END

% >>> CONCLUSION-BEGIN
\section{Conclusion}

Whether AI-based crop-yield research is distributed in proportion to where crops matter, where
evidence is most needed, or where research can most readily be produced is an empirical
question that publication counts cannot answer on their own. The analysis addresses this question by coding the empirical location of each study from its
study-area, data, and methods sections rather than from the affiliations of its authors. Each
study is allocated fractionally across the country-crop systems it examines, and the resulting
research intensities are set against harvested area, production, a nine-component measure of
research need that excludes agricultural scale by design, and national capacity-related
indicators.

Geographic coverage is limited: 122 of 199 countries carry any eligible study, 485 of 2,616
country-crop systems contain research, and wheat, maize, and rice absorb 69.3\% of crop
attention. Harvested area shows the largest and most stable association with research
participation among the covariates examined, and it is the only one whose sign and precision hold
across every specification tried.
Correspondence with agricultural scale is partial, weakening once the crop dimension is
restored. National research and development spending is associated with both participation and
intensity, though the association moves across income groups, and the other capacity-related
indicators show no consistent directional pattern. The need association is too imprecise to
resolve, which limits what can be concluded about need rather than settling it. The mismatch between research standing and need is itself
geographically patterned, and more sharply than the counts it is built from: neighbouring
countries share it, positively across western and northern Europe and negatively across much
of sub-Saharan Africa.

Counting publications by the country of their authors, or counting them at all without a
denominator, cannot separate concentration that follows agriculture from concentration that
departs from it. Placing study-location research attention alongside agricultural scale, research need, and
scientific capacity in one country-crop frame separates them, and a large share of the world's
production systems carries no evidence at all. Nothing in the analysis establishes that models
fitted in well-studied settings perform poorly elsewhere, because cross-country predictive
transfer is not tested here. The uneven geographic coverage, and the large number of country-crop
systems with no resolved evidence, limit the extent to which this literature can be treated as
representative of global cropping systems.

Widening that coverage depends on practices the field can adopt directly. Study locations
stated explicitly would remove the inference now needed to place more than half the corpus.
Validation designs that hold out geography rather than observations would establish whether
models generalise across geographic domains beyond the settings that produced them. Neither practice substitutes for the
capacity constraints that shape where agricultural research can be produced, and those
constraints fall hardest on the country-crop systems where the shortfall relative to need is
largest.
% <<< CONCLUSION-END

\bmsubsection*{Author Contributions}

Prasamsha Poudel and Saugat Khanal contributed equally to this study.

\bmsubsection*{Conflicts of Interest}

The authors declare no conflicts of interest.

\bmsubsection*{Funding}

No funding was received for this study.

\bmsubsection*{Data Availability Statement}

The bibliographic corpus was assembled from the public OpenAlex database. Agricultural production and harvested-area data are from FAOSTAT, research-need components from FAOSTAT food security, macroeconomic, employment, and temperature domains, scientific-capacity indicators from World Bank Open Data, and country boundaries from Natural Earth; all sources are publicly available. The coded study-level dataset, the derived country-crop panel, the analysis code, and the outputs behind every table and figure are available at \url{https://github.com/Saugat2000/ai-crop-yield-research-alignment}.

\bibliography{manuscript1_wiley}

\appendix
\renewcommand{\thetable}{A\arabic{table}}
\renewcommand{\thefigure}{A\arabic{figure}}

\renewcommand{\thesection}{Appendix \Alph{section}}
\renewcommand{\thetable}{A\arabic{table}}

\bmsection{Search, Screening, Deduplication, and Corpus Validation}\applabel{Appendix A}{app:corpus}

Retrieval ran as two passes over OpenAlex, one restricted to journal and preprint document
types and one to conference papers and abstracts, executed as 86 query expressions. All 86
completed. Retrieval returned 102,892 records. Removing duplicate identifiers left 45,845
unique records. Records sharing a DOI, a normalised title, or an author-year signature were
then resolved into single works, yielding 43,543 unique works. Ambiguous groupings were
split rather than merged where the evidence for merging was weak.

Screening applied six eligibility criteria to the title and abstract of all 43,543 works: the
study concerns an agricultural crop; crop yield is the modelled, predicted, or estimated
outcome; the method is genuinely data-driven or machine-learning based; the work reports an
empirical application; the document type is eligible; and the publication date falls from
2000 onward. Process-based crop simulation without a learned component fails the third
criterion.

Screening used an automated rule set rather than human reading, and the
corpus records the decision method. Exclusions distinguish those resting on
affirmative evidence from those resting on the absence of evidence, because the second class can lose an
eligible study irrecoverably. Error rates were estimated from stratified probability samples adjudicated against the same six
criteria. Eighty-five records were adjudicated across three strata: 45 drawn from the largest
exclusion stratum of 8,677 records, and 20 each from two inclusion strata of 4,143 and 2,902
records, with design weights of 99.7, 82.9, and 58.0. The false-inclusion rate is 10.0\% and the
false-exclusion rate 2.2\%. The point estimates satisfy the pre-specified ceilings; the 95\%
interval upper bounds do not, reaching 0.236 and 0.360 for false inclusion against a ceiling of
0.150 and 0.116 for false exclusion against a ceiling of 0.050, because the adjudicated samples
are too small to bound a rate of that size. Three adjudications returned an uncertain verdict.
Only one exclusion reason was sampled, so the false-exclusion estimate does not generalise to the
other exclusion strata. The implied 257 ineligible studies inside the corpus are reported rather
than removed, and the screening thresholds were not altered after the fact. Adjudication was
model-assisted review against the recorded evidence, not independent human reading, and is
reported on that basis (Table~\ref{tab:validation}).

The corpus contains 7,045 eligible studies: 4,842 journal articles, 1,377 conference
papers, 252 preprints, 249 book chapters, 185 dissertations, 117 conference abstracts, and 23
reports, giving 1,494 conference papers and abstracts combined. Abstracts are present for
5,927 studies. A deterministic keyword screen of titles and abstracts finds explicit machine-learning method
terms in 62.6\% of the corpus, classical-regression terms alone in 9.6\%, process-model terms
alone in 2.9\%, and no method keyword in 24.9\%, mostly where no abstract is available. Resolving
the corpus into model families gives deep learning in 29.6\% of records, tree ensembles in
24.4\%, support-vector and kernel methods in 10.4\%, other machine-learning methods in 40.8\%,
classical statistical methods in 17.4\%, and process-based models in 6.4\%; the families overlap
because studies compare several. Rerunning the main analyses on the 4,261-study subset carrying
an explicit machine-learning term changes none of them materially: the all-country Gini moves
from 0.843 to 0.868, the crop Gini from 0.719 to 0.714, the top-three crop share from 69.3\% to
67.7\%, the participation coefficient on log harvested area from 0.713 to 0.707, the intensity
coefficient on research and development expenditure from 0.522 to 0.508, and Moran's $I$ on the
mismatch measure from 0.353 to 0.291, with $p = 0.0001$ throughout. The breadth of the
eligibility definition does not drive the reported geography. The corpus was frozen on 13 August 2026 and is fixed by the SHA-256 checksum
\texttt{98c132b0\ldots26024c}; every dataset used in the analysis derives from that frozen
set.

\bmsection{Study-Location and Crop Coding Protocol}\applabel{Appendix B}{app:coding}

Study location is coded from the study's own account of where its agricultural data
originated. A country name alone is never sufficient evidence. The name must appear in a
context that positively asserts an empirical study area, such as a statement that field
trials were conducted or yield records obtained in a named place. Country names occurring
within an institution name, a funder, a cultivar name, or a citation are evidence about
authorship or provenance rather than study area, and are rejected with the reason recorded.

Each resolved record carries the cue type that produced the assignment. The value
\texttt{study\_area} marks an explicit study-area statement. The value
\texttt{locative\_only} marks an assignment resting on a weaker locative phrase; such
assignments are capped at medium confidence and flagged so that they can be dropped in
sensitivity analysis. A citation-context filter prevents a country named only as the subject
of previously published work from being read as a study area.

Three outcomes are held distinct from a resolved single country. Studies whose scope is
global or continental receive a scope label and no single country code, because assigning a
global study to one country would create an observation the study does not support. Ambiguous
toponyms, including Georgia, Jordan, and Chad, are a country in one reading and a subnational
unit, a person, or a common noun in another; these require a disambiguating cue and otherwise
return an ambiguous result. Studies naming no study area return unresolved, recorded as a
distinct state rather than as an absence of research.

Location coding returns five mutually exclusive states: 1,821 studies resolve to a single
country, 56 to several named countries, 915 are global in scope, 268 regional, and 3,985
unresolved, summing to the 7,045-study corpus; 43.4\% carry some usable location
signal. Of the country-resolved studies, 869 rest on the weaker locative cue. For crops,
4,454 studies name at least one identified crop, 654 more than one, and 2,364 none.

Location and crop coding were produced by an automated, model-assisted extractor. Every record
stores the value, the verbatim evidence text, the cue type, a confidence level, and a
verification status; all 7,045 records are stored as unverified, and none was confirmed by
independent human reading.

Coding accuracy was estimated from a stratified probability sample of 210 records drawn under
seed 20260831 across seven strata spanning study-area and locative cues, multi-country,
global-scope, regional, and unresolved records, with design weights from 2.8 to 184.6. Verdicts
were assigned by model-assisted adjudication against the title and abstract, which is a first
pass requiring human confirmation rather than independent review. Two quantities differ and are
reported separately. Allocation accuracy, meaning that the countries receiving research weight
are the countries the study examines, is 90.4\% (123 of 136 determinable records, 95\% CI 84.3 to
94.3); 0.7\% of records are allocated wholly to the wrong country, and 8.8\% miss or add a
country, almost all in the multi-country stratum. Scope labelling is weaker: the global label is
often triggered by generic wording such as global food security or world population, and 11 of 20
sampled regional records name no study area and should be unresolved. Scope-label errors mostly
leave allocation intact, because allocation runs on accepted country mentions rather than the
scope label. Three reproducible false-positive patterns were confirmed in the corpus: all
six records mentioning New South Wales are coded to the United Kingdom, all three mentioning the
Guinea savanna are coded to Guinea, and one of four records coded to Georgia is the United States
state. Recall is the weaker property: 62.0\% of sampled unresolved records (95\% CI 48.2 to 74.1)
name a study area the coder did not recover, projecting to roughly 1,000 of the 3,985 unresolved
studies. Those unrecovered records track the existing allocation closely by income group and by
mean need percentile, so on observable dimensions the under-resolution does not appear
geographically selective. Crop coding precision is 89.4\% (152 of 170, 95\% CI 83.9 to 93.2).
Table~\ref{tab:validation} reports the design and the error taxonomy.

Allocation uses each study's accepted country evidence: the resolved country for single-country
studies, and otherwise every country the coder accepted as an explicit empirical mention. That rule admits some studies whose overall scope label is global or
unresolved but which name concrete study areas in their titles or abstracts. Under it, 2,031
studies carry both accepted country evidence and a FAOSTAT-matched crop: 1,347
single-country, 47 multi-country, 343 global-scope, and 294 unresolved-scope. Each study
contributes a total weight of one, divided equally across the country-crop combinations it
covers, giving 2,966 study-country-crop records whose weights sum exactly to 2,031. Of that
weight, 10.9 study-equivalents fall on country-crop pairs with no FAOSTAT denominator row
and sit outside the 2,616-cell panel, which carries 2,020.2. The allocation rule does not
drive the results: excluding the global-scope allocations leaves the all-country
concentration Gini at 0.847, and restricting allocation to studies whose scope names
specific countries gives 0.839. Full counts, in which a study contributes one to every cell
it covers, are stored alongside the fractional counts and used only for descriptive
comparison.

\bmsection{Variable Definitions, Sources, Coverage, and Index Construction}\applabel{Appendix C}{app:variables}

\bmsubsection*{Agricultural scale}

The FAOSTAT crop extract covers 1990 to 2024 and contains 111,378 country-crop-year
observations, recording harvested area in hectares, production in tonnes, and yield in tonnes
per hectare. Where FAOSTAT does not report yield directly, yield is derived as production
divided by harvested area, and the derivation is recorded in the source field.

The panel universe is defined as follows. Take the 199 countries and 25 crops carried through
the identifier crosswalk, giving 4,975 possible combinations. Retain the combination where
FAOSTAT reports the crop in that country for any year from 1990 to 2024, which gives 2,590 cells,
and add 30 cells for Guadeloupe, French Guiana, Martinique, and R\'eunion, which the analytical
layer holds separately but FAOSTAT reports inside France. Four cells present in FAOSTAT
(Jordan-pulses, Chad-yam, Saint Vincent and the Grenadines-sugarcane, Yemen-sugarcane) fall
outside the panel through the crosswalk. The result is 2,616 cells. Of these, 109 carry no
FAOSTAT area, production, or yield value for the reference period and support no
denominator-based comparison; none of the 109 carries research, so restricting the panel to the
2,507 cells with an observation moves the share of studied systems from 18.5\% to 19.3\%.

Each panel cell takes the mean of the five reference years 2018 to 2022. The number of reference
years actually observed is stored with each cell: 2,464 of the 2,616 cells draw on all five. Production values are present for 2,505 cells, of which 2,345 are strictly
positive; harvested area is present for 2,504 cells, of which 2,343 are positive. Cells
carrying a recorded zero are retained and stay distinct from cells with no observation.

Country identifiers are harmonised through a crosswalk linking FAOSTAT M49 codes, ISO
alpha-3 codes, World Bank names, and the alpha-2 codes OpenAlex uses for affiliations. An
unmatched identifier is logged rather than dropped. FAOSTAT aggregate areas, including World
and continental groupings, are flagged and excluded from country-level analysis.

Crops are standardised to 25 items spanning cereals, roots and tubers, pulses, oilseeds,
sugar crops, and the fibre and beverage crops appearing in the corpus. Several items group
closely related FAOSTAT commodities that the study literature does not distinguish, including
banana and plantain, oat and rye, and coffee and cocoa. Crops named in eligible studies but
lacking a separate FAOSTAT item, teff among them, cannot receive a production denominator and
are recorded as such rather than folded into a residual category.

\bmsubsection*{Research-need components and index construction}

The research-need index combines the nine non-scale indicators listed in
Table~\ref{tab:vardetail}; the two agricultural-scale shares are recorded with the need
data but excluded from the index because they are the paper's importance measures.
Indicators are recorded with their declared direction, converted to percentile ranks, and
combined into rank-aggregation, equal-weighting, entropy-weighting, and first-principal-
component indices, plus the eleven-component variant used for sensitivity. Each index is also stored as
a percentile on a zero-to-one scale. Component coverage varies and is stored for every country. The floor is five of the nine primary
components, and it operates in three cases: below the floor no index is computed; at or above it
but incomplete, the index is the mean of the components observed; complete, the index uses all
nine. Applying the floor to the nine primary components indexes 188 of the 199 countries. The
published index applied the same floor to the eleven recorded indicators, which include the two
agricultural-scale shares that are never missing, and indexed 193 countries, five of
which observe fewer than five actual need components. The correction leaves the reported results
unchanged, moving the participation coefficient on need from 0.944 to 0.950 and the mismatch
clustering statistic from 0.363 to 0.346. The nine components are weakly related rather than
interchangeable: Cronbach's alpha over the harmonised ranks is 0.669, the mean pairwise Spearman correlation is 0.243 (Figure~\ref{fig:needcorr}), and the first
principal component explains 41.9\% of variance, so the index
summarises several distinct dimensions rather than one latent construct.

\bmsubsection*{Scientific-capacity indicators}

Scientific capacity uses six World Bank indicators retrieved through the World Bank API and
verified against the API's own record counts. Each country takes the most recent non-missing
observation, and the year of that observation is stored beside the value so that staleness
remains visible. Population required repair: the initial retrieval was truncated by an API
paging limit, losing 1,540 observations across 42 entities, and the series was re-fetched and
verified before use. All other indicators reconciled exactly against the API counts. Missing
capacity values are not imputed, so models including research and development expenditure
estimate on a smaller sample than models excluding it.

\bmsubsection*{Spatial layer}

Country geometries are the Natural Earth admin-0 boundaries in geographic coordinates on the
WGS 84 datum (EPSG:4326). The analytical layer holds 195 countries, 93 represented as single
polygons and 102 as multipolygons. Each record carries the FAOSTAT M49 code, ISO alpha-3 code, World Bank name, region, and income
group, together with a polygon centroid and a land area computed on the Equal Earth projection.
Nearest neighbours are found among those polygon centroids after projection to Equal Earth, using
planar Euclidean distance, and every matrix is row standardised. Equal Earth preserves area
rather than distance, so the statistics were recomputed with neighbours found on true
great-circle distance: the neighbour set changes for 86 of 195 countries, and Moran's $I$ on the
mismatch measure is 0.364 against 0.359 at $k=6$, 0.373 against 0.364 at $k=4$, and 0.350 against
0.333 at $k=8$, with local clusters of 16 High--High and 5 Low--Low against 17 and 6. Aggregate and historical entities are excluded
and the reason for each exclusion is recorded. Four panel territories have no matching
boundary; all four carry no eligible study.



\bmsection{Robustness and Sensitivity Specifications}\applabel{Appendix D}{app:robust}

Table~\ref{tab:specs} lists the specifications behind Section~4.6. Concentration variants
re-compute the distribution of research attention across countries; model variants re-estimate
the participation specification of Equation~\eqref{eq:logit}; spatial variants re-compute
Equation~\eqref{eq:moran}.

Two estimation limitations are recorded here rather than in the body. The negative binomial
model of intensity converged but did not return standard errors under the clustered covariance
estimator, so it informs the overdispersion question only and no interval is read from it. A
hurdle specification separating participation from positive counts was fitted during model
development, but its coefficients were not retained in the reported output and it is not used.

\begin{table*}[!t]
\centering
\scriptsize
\setlength{\tabcolsep}{3.5pt}
\renewcommand{\arraystretch}{1.1}
\caption{Robustness and sensitivity specifications}
\label{tab:specs}
\begin{tabular}{@{}>{\raggedright\arraybackslash}p{2.9cm}
                  >{\raggedright\arraybackslash}p{2.6cm}
                  >{\raggedright\arraybackslash}p{3.0cm}
                  >{\raggedright\arraybackslash}p{3.1cm}
                  >{\raggedright\arraybackslash}p{2.3cm}@{}}
\toprule
Check & Main specification & Alternative & Purpose & Sample affected \\
\midrule
Counting rule & Fractional, Eq.~\eqref{eq:frac} & Full counting &
Test whether concentration depends on how multi-country studies are divided & All studies \\
\addlinespace
Attention weighting & Unweighted study count & Citation-weighted count &
Test whether weighting by attention received changes the distribution & All studies \\
\addlinespace
Document type & Journal and conference & Journal only; preprints excluded &
Test whether conference and preprint literature drives the pattern & 1,712 studies journal-only \\
\addlinespace
Location confidence & All resolved locations & Locative-cue cells excluded; low-confidence excluded &
Test whether weaker location coding drives the pattern & Cells resolved only by locative cue \\
\addlinespace
Influential countries & All countries & Exclude USA; China; both; and with India &
Test whether the largest contributors produce the concentration & Up to three countries \\
\addlinespace
Period & Full 2000--2026 & 2000--2014; 2015--2019; 2020--2026 &
Test whether the pattern is specific to one period & All studies \\
\addlinespace
Screening error & Frozen corpus & Adjustment for the measured false-inclusion rate &
Bound the effect of screening error on concentration & Whole corpus \\
\addlinespace
Need index & Rank aggregation & Equal weighting; entropy weighting; first principal component &
Test whether the need result depends on aggregation & All indexed countries \\
\addlinespace
Income group & Pooled & Low, lower-middle, upper-middle, high separately &
Test whether associations differ by income & Estimation sample split \\
\addlinespace
Spatial weights & kNN, $k=6$ & $k=4$; $k=8$; queen; distance band &
Test whether spatial results depend on the neighbour definition & Queen drops 43 countries \\
Neighbour distance & Equal Earth planar & Great-circle distance &
Test whether an equal-area projection distorts the neighbour sets & 86 of 195 neighbour sets change \\
Rank ties & Average ranks & Minimum; maximum; dense; standardised log count &
Test whether mismatch clustering depends on the tie convention & 71 countries tied at zero \\
Need-index coverage floor & Five of eleven recorded indicators & Five, seven, or nine of the nine primary components &
Test whether thinly measured countries drive the need result & 188 countries at the corrected floor \\
Area elasticity & Fixed at one by the offset & Freely estimated, with a Wald test &
Test the exposure restriction in the intensity model & 1,792 cells \\
Method definition & Full eligible corpus & Studies naming an explicit machine-learning method &
Test whether eligibility breadth drives the geography & 4,261 studies \\
\bottomrule
\end{tabular}
\begin{tablenotes}
\item[] Sample affected refers to the units whose inclusion or weighting changes relative to the main specification.
\end{tablenotes}
\end{table*}


\begin{table*}[!t]
\centering
\footnotesize
\caption{Validation of screening, study-location coding, and crop coding}
\label{tab:validation}
\begin{tabular*}{\textwidth}{@{\extracolsep\fill}l r r r r r l@{}}
\toprule
\multicolumn{7}{l}{\textit{Panel A. Study-location coding by stratum}} \\
Stratum & Records & Sampled & Weight & Correct & Rate & 95\% CI \\
\midrule
Global scope, countries accepted & 343 & 30 & 11.4 & 0 & 0.000 & [0.000, 0.114] \\
Multiple named countries & 56 & 20 & 2.8 & 8 & 0.400 & [0.219, 0.613] \\
Regional scope & 268 & 20 & 13.4 & 5 & 0.250 & [0.112, 0.469] \\
Single country, locative cue & 859 & 45 & 19.1 & 34 & 0.756 & [0.613, 0.858] \\
Single country, study-area cue & 962 & 45 & 21.4 & 43 & 0.956 & [0.852, 0.988] \\
Unresolved scope, countries accepted & 294 & 30 & 9.8 & 2 & 0.067 & [0.018, 0.213] \\
Unresolved, not allocated & 3,691 & 20 & 184.6 & 11 & 0.550 & [0.342, 0.742] \\
\midrule
\multicolumn{7}{l}{\textit{Panel B. Main quantities}} \\
\multicolumn{4}{@{}l}{Location coding accuracy, allocated records (drive the counts)} & 87 & 0.551 & [0.473, 0.626] \\
\multicolumn{4}{@{}l}{Location coding accuracy, single-country study-area cue} & 43 & 0.977 & [0.882, 0.996] \\
\multicolumn{4}{@{}l}{Location coding accuracy, single-country locative cue} & 34 & 0.971 & [0.855, 0.995] \\
\multicolumn{4}{@{}l}{Country-assignment precision, allocated records} & 87 & 0.554 & [0.476, 0.630] \\
\multicolumn{4}{@{}l}{False-unresolved rate (study area present but not coded)} & 31 & 0.620 & [0.482, 0.741] \\
\multicolumn{4}{@{}l}{Crop coding precision (records with a determinable crop)} & 152 & 0.894 & [0.839, 0.932] \\
\midrule
\multicolumn{7}{l}{\textit{Panel C. Screening validation}} \\
Largest exclusion stratum & 8,677 & 45 & 99.7 & 44 & 0.978 & [0.884, 0.996] \\
Inclusion, all criteria affirmative & 4,143 & 20 & 82.9 & 19 & 0.950 & [0.764, 0.992] \\
Inclusion, empirical marker not named & 2,902 & 20 & 58.0 & 14 & 0.700 & [0.481, 0.855] \\
\bottomrule
\end{tabular*}
\begin{tablenotes}
\item[] Panel A reports study-location coding accuracy by stratum on the 210-record probability
sample (seed 20260831); Weight is the design weight, and Rate counts a record correct only if
both the scope label and the country assignment match the title and abstract. That strict rule
explains the zero rate for the global-scope stratum: the scope label was triggered by generic
wording such as global food security in every sampled record, while the countries those studies
actually name were still allocated correctly. Scope-label errors of this kind do not misdirect
research weight, because allocation runs on accepted country mentions rather than on the scope
label, which is why Panel B reports allocation accuracy separately. Panel B reports the quantities that bear on the analysis, computed on records whose verdict is determinable. Panel C reports the screening validation of 85 records. Intervals are Wilson intervals. All adjudication was model-assisted review against the recorded evidence rather than independent human reading, and requires human confirmation before it is treated as settled.
\end{tablenotes}
\end{table*}





\begin{table*}[!t]
\centering
\caption{Specification checks on the participation and intensity models}
\label{tab:specchecks}
\begin{tabular*}{\textwidth}{@{\extracolsep\fill}l r r r r r@{}}
\toprule
& Baseline & Crop FE & No R\&D, & Intensity & Intensity, \\
Covariate & participation & participation & expanded & (offset) & free area \\
\midrule
Log harvested area & 0.713 & 0.656 & 0.670 & \multicolumn{1}{c}{---} & 0.849 \\
R\&D expenditure (\% GDP) & 0.315 & 0.300 & \multicolumn{1}{c}{---} & 0.522 & 0.452 \\
Research-need index & 0.944 & 0.682 & 0.692 & 1.261 & 0.950 \\
Log GDP per capita & 0.347 & 0.336 & 0.628 & $-$0.052 & 0.021 \\
Tertiary enrolment & 0.257 & 0.156 & 0.552 & $-$0.985 & $-$0.865 \\
Internet use & 1.312 & 1.037 & 0.321 & 2.179 & 1.845 \\
Log population & 0.111 & 0.308 & 0.184 & $-$0.108 & 0.015 \\
\midrule
Cells & 1,799 & 1,799 & 2,173 & 1,792 & 1,792 \\
\bottomrule
\end{tabular*}
\begin{tablenotes}
\item[] Coefficients only; intervals are in the replication outputs. Crop fixed effects identify the harvested-area association from variation across countries within a crop. The expanded column drops research and development expenditure, which restores 36 countries and 374 cells. The final column estimates the harvested-area elasticity rather than fixing it at one: the estimate is 0.849 (95\% CI [0.754, 0.944]), and a Wald test rejects unit elasticity ($\chi^2_1 = 9.70$, $p = 0.0018$). Standard errors are clustered by country throughout.
\end{tablenotes}
\end{table*}

\begin{table*}[!t]
\centering
\caption{Nested block models of research participation}
\label{tab:blocks}
\begin{tabular*}{\textwidth}{@{\extracolsep\fill}l r r r r r@{}}
\toprule
Model & Covariates & Log-likelihood & McFadden $R^2$ & AIC & BIC \\
\midrule
Constant only & 0 & $-$1008.6 & 0.000 & 2019.2 & 2024.7 \\
Agricultural scale & 1 & $-$704.1 & 0.302 & 1412.2 & 1423.2 \\
Scale $+$ research need & 2 & $-$662.3 & 0.343 & 1330.6 & 1347.1 \\
Scale $+$ capacity indicators & 6 & $-$641.2 & 0.364 & 1296.5 & 1334.9 \\
Scale $+$ need $+$ capacity & 7 & $-$639.8 & 0.366 & 1295.6 & 1339.6 \\
\midrule
\multicolumn{6}{l}{\textit{Likelihood-ratio tests}} \\
Comparison & $\chi^2$ & df & $p$ & $\Delta$ McFadden $R^2$ & \\
Need added to scale & 83.60 & 1 & 0.0000 & 0.0414 & \\
Capacity added to scale & 125.74 & 5 & 0.0000 & 0.0623 & \\
Capacity added to scale+need & 44.97 & 5 & 0.0000 & 0.0223 & \\
Need added to scale+capacity & 2.83 & 1 & 0.0926 & 0.0014 & \\
\bottomrule
\end{tabular*}
\begin{tablenotes}
\item[] Participation logit on the 1,799-cell estimation sample, standard errors clustered by country. Capacity indicators are research and development expenditure, log GDP per capita, tertiary enrolment, internet use, and log population. Differences in fit describe how much each block adds to in-sample explanation and are not evidence of causal importance. The corresponding comparison for the intensity model, scaled by the estimated dispersion, distinguishes neither block from the other.
\end{tablenotes}
\end{table*}

\begin{table*}[!t]
\centering
\caption{Country-level models of research output}
\label{tab:countrylevel}
\begin{tabular*}{\textwidth}{@{\extracolsep\fill}l r@{\hspace{6pt}}l r@{\hspace{6pt}}l@{}}
\toprule
& \multicolumn{2}{c}{Log study count} & \multicolumn{2}{c}{Studies per hectare (PPML)} \\
\cmidrule(lr){2-3}\cmidrule(lr){4-5}
Covariate & Estimate & 95\% CI & Estimate & 95\% CI \\
\midrule
Log harvested area & 0.138 & [0.07, 0.20] & $-$0.473 & [$-$0.63, $-$0.32] \\
R\&D expenditure (\% GDP) & 0.019 & [$-$0.14, 0.18] & 0.278 & [0.07, 0.49] \\
Research-need index & $-$1.083 & [$-$2.07, $-$0.09] & 0.533 & [$-$0.64, 1.70] \\
Log GDP per capita & $-$0.020 & [$-$0.34, 0.30] & 0.244 & [$-$0.22, 0.71] \\
Tertiary enrolment & $-$0.172 & [$-$0.76, 0.41] & $-$0.456 & [$-$1.35, 0.44] \\
Internet use & 1.078 & [0.02, 2.13] & 0.704 & [$-$0.94, 2.35] \\
Log population & 0.450 & [0.33, 0.57] & 0.363 & [0.18, 0.55] \\
\midrule
Countries & \multicolumn{2}{c}{142} & \multicolumn{2}{c}{142} \\
\bottomrule
\end{tabular*}
\begin{tablenotes}
\item[] Countries are the unit. The left column regresses the log of one plus the country's fractional study count on the same covariates as Table~\ref{tab:reg}; the right column is Poisson pseudo-maximum likelihood with total harvested area as an exposure offset. Heteroskedasticity-robust standard errors. Estimates describe association, not effect.
\end{tablenotes}
\end{table*}

\begin{table*}[!t]
\centering
\footnotesize
\caption{Resolved versus unresolved studies}
\label{tab:resolution}
\begin{tabular*}{\textwidth}{@{\extracolsep\fill}l r r r r r r@{}}
\toprule
& \multicolumn{2}{c}{Resolved (2,031)} & \multicolumn{2}{c}{Unresolved (5,014)} & & \\
\cmidrule(lr){2-3}\cmidrule(lr){4-5}
Measure & Mean & Median & Mean & Median & Statistic & $p$ \\
\midrule
Publication year & 2020.6 & 2022 & 2021.1 & 2023 & 4,567,372 & $<$0.001 \\
Citation count & 23.0 & 4 & 15.4 & 2 & 5,811,177 & $<$0.001 \\
\midrule
\multicolumn{7}{l}{\textit{Percentages}} \\
Journal article & \multicolumn{2}{c}{75.9} & \multicolumn{2}{c}{65.8} & 68.2 & $<$0.001 \\
Conference paper or abstract & \multicolumn{2}{c}{15.7} & \multicolumn{2}{c}{23.4} & 51.2 & $<$0.001 \\
Preprint & \multicolumn{2}{c}{3.3} & \multicolumn{2}{c}{3.7} & 0.3 & 0.557 \\
Other document type & \multicolumn{2}{c}{5.0} & \multicolumn{2}{c}{7.1} & 9.8 & 0.002 \\
Abstract available & \multicolumn{2}{c}{93.9} & \multicolumn{2}{c}{80.2} & 204.8 & $<$0.001 \\
Explicit machine-learning method term & \multicolumn{2}{c}{49.8} & \multicolumn{2}{c}{62.4} & 93.2 & $<$0.001 \\
\bottomrule
\end{tabular*}
\begin{tablenotes}
\item[] Resolved studies are the 2,031 carrying accepted country evidence and a FAOSTAT-matched crop, which enter the allocation; the remaining 5,014 do not. Continuous measures are compared by Mann--Whitney $U$ and percentages by chi-square on the two-by-two table. With samples of this size, small differences reach conventional significance: the publication-year difference is half a year, whereas the citation and abstract-availability differences are larger. The comparison characterises observable selection and does not establish that non-resolution is geographically random.
\end{tablenotes}
\end{table*}

\begin{table*}[!t]
\centering
\footnotesize
\caption{Covariate vintage, dispersion, and spatial diagnostics}
\label{tab:diagnostics}
\begin{tabular*}{\textwidth}{@{\extracolsep\fill}l r r r r r r@{}}
\toprule
\multicolumn{7}{l}{\textit{Panel A. Reference year of the R\&D expenditure observation, 142 estimation countries}} \\
Income group & Countries & Median & Earliest & Latest & \multicolumn{2}{c}{} \\
\midrule
High income & 54 & 2023 & 2015 & 2024 & \multicolumn{2}{c}{} \\
Low income & 14 & 2018 & 2005 & 2024 & \multicolumn{2}{c}{} \\
Lower middle income & 31 & 2019 & 2002 & 2024 & \multicolumn{2}{c}{} \\
Upper middle income & 43 & 2023 & 1999 & 2024 & \multicolumn{2}{c}{} \\
\midrule
\multicolumn{7}{l}{\textit{Panel B. Intensity model: clustered against quasi-Poisson standard errors}} \\
Covariate & Estimate & SE (clustered) & SE (quasi-Poisson) & Ratio & $p$ (clust.) & $p$ (quasi) \\
\midrule
R\&D expenditure (\% GDP) & 0.522 & 0.123 & 0.601 & 4.9 & $<$0.001 & 0.385 \\
Research-need index & 1.261 & 0.590 & 3.637 & 6.2 & 0.033 & 0.729 \\
Log GDP per capita & $-$0.052 & 0.238 & 1.390 & 5.8 & 0.827 & 0.970 \\
Tertiary enrolment & $-$0.985 & 0.425 & 2.892 & 6.8 & 0.020 & 0.733 \\
Internet use & 2.179 & 0.835 & 5.039 & 6.0 & 0.009 & 0.665 \\
Log population & $-$0.108 & 0.056 & 0.291 & 5.2 & 0.054 & 0.711 \\
\midrule
\multicolumn{7}{l}{\textit{Panel C. Lagrange Multiplier tests on country-mean residuals, $k = 6$}} \\
& \multicolumn{2}{c}{Participation logit} & \multicolumn{2}{c}{Intensity PPML} & \multicolumn{2}{c}{} \\
\cmidrule(lr){2-3}\cmidrule(lr){4-5}
Test & Statistic & $p$ & Statistic & $p$ & \multicolumn{2}{c}{} \\
\midrule
LM lag & 0.674 & 0.412 & 0.000 & 0.986 & \\
Robust LM lag & 0.628 & 0.428 & 0.545 & 0.460 & \\
LM error & 0.567 & 0.451 & 0.242 & 0.623 & \\
Robust LM error & 0.521 & 0.470 & 0.787 & 0.375 & \\
\bottomrule
\end{tabular*}
\begin{tablenotes}
\item[] Panel A gives the year of the most recent research and development expenditure observation for each estimation country. Panel B compares the reported cluster-robust standard errors with quasi-Poisson standard errors formed from the Pearson dispersion, $\hat\phi_{P} = 288$. The deviance-based dispersion for the same fit is $\hat\phi_{D} = 1.4$. The two disagree by more than two orders of magnitude because the Pearson statistic is dominated by a small number of cells with very large counts, so a constant variance-to-mean ratio does not describe these data and the quasi-Poisson interval is conservative rather than correct. Cluster-robust errors, which assume no variance function, remain the reported choice. Panel C reports Lagrange Multiplier tests for a spatial lag and a spatial error process on country-mean residuals under the primary six-nearest-neighbour matrix; none rejects at the 5\% level, which agrees with the residual Moran statistics in Section~5.4.
\end{tablenotes}
\end{table*}

\begin{table*}[!t]
\centering
\footnotesize
\caption{Participation models by publication period}
\label{tab:periods}
\begin{tabular*}{\textwidth}{@{\extracolsep\fill}l r@{\hspace{6pt}}l r@{\hspace{6pt}}l@{}}
\toprule
& \multicolumn{2}{c}{Published 2000--2019} & \multicolumn{2}{c}{Published 2020--2026} \\
\cmidrule(lr){2-3}\cmidrule(lr){4-5}
Covariate & Estimate & 95\% CI & Estimate & 95\% CI \\
\midrule
Log harvested area & 0.750 & [0.60, 0.90] & 0.748 & [0.64, 0.86] \\
R\&D expenditure (\% GDP) & 0.567 & [0.26, 0.88] & 0.285 & [0.03, 0.54] \\
Research-need index & 0.617 & [$-$1.92, 3.16] & 0.447 & [$-$1.38, 2.28] \\
Log GDP per capita & 0.073 & [$-$0.75, 0.90] & 0.273 & [$-$0.42, 0.97] \\
Tertiary enrolment & $-$0.767 & [$-$2.21, 0.68] & 0.419 & [$-$0.53, 1.36] \\
Internet use & 1.644 & [$-$0.87, 4.16] & 1.234 & [$-$0.93, 3.40] \\
Log population & $-$0.031 & [$-$0.23, 0.17] & 0.070 & [$-$0.16, 0.30] \\
\midrule
Country-crop cells & \multicolumn{2}{c}{1,799} & \multicolumn{2}{c}{1,799} \\
Cells carrying a study & \multicolumn{2}{c}{225} & \multicolumn{2}{c}{402} \\
\bottomrule
\end{tabular*}
\begin{tablenotes}
\item[] The outcome is whether a country-crop system carries a study published in the stated period; the covariate panel is unchanged, because the capacity and need indicators are single most-recent observations rather than annual series. The models test whether the associations are stable across the literature's recent expansion, not whether covariates changed. Standard errors are clustered by country. Note that 2026 is a partial publication year in the corpus.
\end{tablenotes}
\end{table*}

\bmsection{Supplementary Figures}\applabel{Appendix E}{app:figs}

Figure~\ref{fig:crops} accompanies Section~5.1, Figure~\ref{fig:coef} accompanies
Section~5.4, where the corresponding estimates are reported in Table~\ref{tab:reg}, and
Figures~\ref{fig:temporalgini} and~\ref{fig:needcorr} accompany Sections~5.6 and~4.3.

\begin{figure*}[!t]
\centering
\includegraphics[width=0.98\linewidth,height=0.68\textheight,keepaspectratio]{fig_02_v2_crop_attention_area.pdf}
\caption{Research attention and world harvested-area share, by crop, for the fifteen
crops with the largest share of fractional research attention. Harvested area is the
cross-crop benchmark because hectares, unlike tonnes, are comparable across crops.}
\label{fig:crops}
\end{figure*}

\begin{figure*}[!t]
\centering
\includegraphics[width=0.98\linewidth,height=0.68\textheight,keepaspectratio]{fig_06_v2_coefficients.pdf}
\caption{Estimated coefficients for the participation and intensity models. Bars give 95\% confidence intervals. The estimates are those reported in Table~\ref{tab:reg}.}
\label{fig:coef}
\end{figure*}

\clearpage





\begin{figure*}[!t]
\centering
\includegraphics[width=0.80\linewidth,height=0.46\textheight,keepaspectratio]{fig_A4_temporal_concentration.pdf}
\caption{Gini coefficient of country research counts by publication period, computed across all
195 countries in the boundary layer. The vertical axis spans 0.50 to 1.00 so that the small
movement between periods is not visually inflated, and the dashed line marks the full-period
value of 0.848. The 2026 point is drawn separately because that publication year is incomplete
in the corpus.}
\label{fig:temporalgini}
\end{figure*}

\begin{figure*}[!t]
\centering
\includegraphics[width=0.78\linewidth,height=0.58\textheight,keepaspectratio]{fig_A5_need_component_correlations.pdf}
\caption{Spearman correlations among the nine research-need components, after direction
harmonisation and conversion to percentile ranks. The components are weakly related rather than
interchangeable.}
\label{fig:needcorr}
\end{figure*}

\clearpage
\end{document}
